\section{The completed three-\texorpdfstring{$Z$}{Z} master and common-prefix geometry}
\label{sec:threeZ}

The three principal $Z$-faces are not three independent sources. We use the
completed three-face Cauchy master and diagonal Mellin representation of
Heap--Li--Zhao~\cite{HLZ}. At fixed height the three faces arise from one
completed shifted product and one joint three-variable residue. Only after
this completion is the diagonal Mellin variable introduced. In the general
HLZ twisted formula the auxiliary outer Dirichlet polynomial produces indices
$h,k$ and the diagonal relation $hm=kn$. For $(h,k)=1$ this is
\begin{equation}
\boxed{m=k\ell,\qquad n=h\ell.}
\label{eq:hlz-diagonal}
\end{equation}
The invariant contour form of this paper, however, contains no such auxiliary
outer Dirichlet polynomial: its additional test factor is the packet
polarization $\mathcal B_{v,w}$, which is retained as a smooth holomorphic
height/Mellin weight. Thus the actual source specialization used below is
$h=k=1$, and \eqref{eq:hlz-diagonal} reduces to $m=n=\ell$.

\begin{proposition}[Trivial outer twist of the invariant packet source]
\label{prop:trivial-outer-twist}
In the application of the HLZ diagonal calculation to the contour matrix
$A_T$, the auxiliary outer Dirichlet polynomial is identically one. Hence its
outer indices satisfy
\[
\boxed{g=h=k=1.}
\]
If this source is embedded in any finite auxiliary divisor range $d\le Y$, its
Hilbert-valued coefficient family has only the coordinate $d=1$ nonzero.
Consequently the outer packet translations are the identity, the common
divisor map of Proposition~\ref{prop:common-divisor-feature} restricts to the
identity, and the factor $de/g$ equals $1$. In particular the outer excess
annulus \eqref{eq:outer-excess-annulus} is empty for the invariant contour
source.
\end{proposition}

\begin{proof}
The exact same-form identity \eqref{eq:same-form-identity} contains the shifted
Hardy product and the packet polarization, but no factor of the form
$Q_{\rm out}(s)\overline{Q_{\rm out}}(1-s)$. The HLP series in
\eqref{eq:Hex-variable} belongs to the logarithmic-derivative source and is not
an auxiliary outer test polynomial. In the HLZ twisted formula the indices
$g,h,k$ arise only after expanding such an outer Dirichlet polynomial and its
adjoint. Therefore the present contour form is the specialization
$Q_{\rm out}\equiv1$, giving $g=h=k=1$. The packet polarization remains in the
analytic weight before the finite shift residues, so no discrete outer index is
created by it. This specialization removes only the auxiliary outer indices
$g,h,k$; it does not identify or contract the internal HLP radial carrier $M$.
In the unconditional decomposition used in the final proof that carrier remains
inside the direct one-$\chi$ exact--model residual and is handled by the shell
sideband decomposition. The remaining assertions follow immediately from the
definitions of the translation and common-divisor maps.
\end{proof}

\begin{remark}[Scope of the scalar HLZ input]
We use the completed three-face Cauchy identity and diagonal calculation of
Heap--Li--Zhao~\cite{HLZ}, specialized as in
Proposition~\ref{prop:trivial-outer-twist}.  Their scalar lemma chooses the
auxiliary function after the Cauchy variables have been fixed; it does not
provide a single function analytic jointly in the Mellin variable and in all
shift variables through the collision locus.  We therefore never impose such a
joint analyticity requirement.  The packet polarization is not identified with
their auxiliary Dirichlet polynomial.  For the finite model source used in the
final proof, the passage from this scalar identity to the packet matrix is made
below by the exact one-sided packet factorization and Fubini.  A stronger
packet-bilinear transfer for the exact low-level HLP completion is recorded
later only as an optional comparison route and is not used in the proof of
Theorem~\ref{thm:main}.
\end{remark}

Fix constants
\[
0<r_1<r_2<r_3,
\qquad
r_{j+1}-r_j\ge r_*>0,
\]
large enough that the three circles
\begin{equation}
\mathcal T_{j,L}:=\{z:|z|=r_j/L\}
\label{eq:HLZ-separated-circles}
\end{equation}
enclose every retained physical shift.  Their product
$\mathcal T_L:=\mathcal T_{1,L}\times\mathcal T_{2,L}\times\mathcal T_{3,L}$
is a separated Cauchy torus: for $\mathbf z\in\mathcal T_L$,
\begin{equation}
|z_i-z_j|\ge r_*/L\qquad(i\ne j).
\label{eq:HLZ-shift-separation}
\end{equation}
For such $\mathbf z$ define
\begin{equation}
\boxed{
H_{\mathbf z}(u):=
\exp(u^2)
\frac{(2u-z_2+z_1)(2u-z_2+z_3)}
     {(z_1-z_2)(z_3-z_2)}.
}
\label{eq:HLZ-explicit-auxiliary}
\end{equation}
Then, pointwise on $\mathcal T_L$,
\begin{equation}
H_{\mathbf z}(0)=1,
\qquad
H_{\mathbf z}\!\left(\frac{z_2-z_1}{2}\right)=0,
\qquad
H_{\mathbf z}\!\left(\frac{z_2-z_3}{2}\right)=0.
\label{eq:HLZ-auxiliary-properties}
\end{equation}
There is no contradiction at $\mathbf z=0$, because
$\mathbf z=0\notin\mathcal T_L$ and no jointly analytic extension across the
collision locus is asserted.  Moreover on every vertical line used below,
$H_{\mathbf z}(u)$ is a polynomial factor in $L(1+|u|)$ times the Gaussian
$e^{-\Im(u)^2}$, uniformly on the separated torus.

\begin{proposition}[Exact scalar three-face Cauchy identity]
\label{prop:HLZ-three-face-Cauchy}
Let $X>0$ and write
\[
\mathfrak Z_{\alpha,\beta,\gamma}(X)
:=Z_{\alpha,\beta,\gamma,1,1}
+X^{-\alpha-\beta}Z_{-\beta,-\alpha,\gamma,1,1}
+X^{-\beta-\gamma}Z_{\alpha,-\gamma,-\beta,1,1}.
\]
With $\Delta(\mathbf z)=\prod_{i<j}(z_j-z_i)$, the residue identity used by
Heap--Li--Zhao is, after the trivial outer-twist specialization,
\begin{equation}
\boxed{
\begin{aligned}
\mathfrak Z_{\alpha,\beta,\gamma}(X)
={}&-\frac12X^{-(\alpha+\beta+\gamma)/2}
\frac1{(2\pi i)^3}
\iiint_{\mathcal T_L}
Z_{z_1,-z_2,z_3,1,1}\,
\frac{\Delta(\mathbf z)^2
X^{(z_1-z_2+z_3)/2}}
{\prod_{j=1}^3(z_j-\alpha)(z_j+\beta)(z_j-\gamma)}
\,d\mathbf z .
\end{aligned}}
\label{eq:HLZ-three-face-Cauchy}
\end{equation}
The three contours are chosen from the outset with the distinct radii
\eqref{eq:HLZ-separated-circles}, small enough for the local residue formula
and large enough to enclose the retained physical shifts; they are also chosen
so that none of the polar hypersurfaces of
$Z_{z_1,-z_2,z_3,1,1}$ meets the integration torus.
\end{proposition}

\begin{proof}
This is formula~(28) in Heap--Li--Zhao~\cite{HLZ}, with $h=k=1$.
Their statement allows small circles of radii $\ll1/L$ enclosing the physical
shifts.  Because only finitely many local polar hypersurfaces occur, the radii
may be chosen distinct and generic while retaining the same enclosed Cauchy
poles; the residue identity is then evaluated on that fixed torus.  We do not
assert a deformation through a collision hypersurface.  The separated choice
is important because it permits the auxiliary function
\eqref{eq:HLZ-explicit-auxiliary} to be chosen pointwise in $\mathbf z$ without
requiring a jointly analytic extension across $z_i=z_j$.
\end{proof}

\begin{lemma}[Parametric form of the HLZ diagonal contour]
\label{lem:parametric-HLZ-diagonal}
Put $\log_3T:=\log\log\log T$. Write
\[
 Z_{\mathbf z,1,1}:=Z_{z_1,-z_2,z_3,1,1},
 \qquad
 D_{\mathbf z}(s):=
 \frac{\zeta(1+z_1-z_2+2s)\zeta(1-z_2+z_3+2s)}
      {\zeta(1-z_2+2s)}.
\]
Fix a constant $c_+>0$, large enough in terms of the radii in
\eqref{eq:HLZ-separated-circles}, so that $D_{\mathbf z}(s)$ is represented by
its absolutely convergent Euler product on $\Re s=c_+/L$.  Let
$K_{T,\mathbf z}(s)$ be analytic on the rectangle between
$\Re s=c_+/L$ and $\Re s=-1/\log_3T$, with
\[
 K_{T,\mathbf z}\!\left(\frac{z_2-z_1}{2}\right)=
 K_{T,\mathbf z}\!\left(\frac{z_2-z_3}{2}\right)=0.
\]
Assume that for some fixed $C_K,c_K>0$, uniformly on that strip,
\begin{equation}
 \left|\mathcal D K_{T,\mathbf z}(\sigma+iv)\right|
 \ll L^{C_K}(1+|v|)^{C_K}e^{-c_Kv^2}
 \label{eq:parametric-HLZ-K-bound}
\end{equation}
for every operator $\mathcal D$ in a prescribed finite family of packet
parameter derivatives, including $\mathcal D=1$.  Define the diagonal Mellin
integral on the absolutely convergent right line by
\begin{equation}
 \widetilde Z_{\mathbf z,1,1}[K]
 :=\frac1{2\pi i}\int_{c_+/L-i\infty}^{c_+/L+i\infty}
 \frac{K_{T,\mathbf z}(s)}s T^{2s}D_{\mathbf z}(s)\,ds.
 \label{eq:parametric-HLZ-diagonal-integral}
\end{equation}
Then, pointwise for $\mathbf z\in\mathcal T_L$,
\begin{equation}
 K_{T,\mathbf z}(0)Z_{\mathbf z,1,1}
 =\widetilde Z_{\mathbf z,1,1}[K]
 +\mathcal E_K(\mathbf z),
 \label{eq:parametric-HLZ-diagonal}
\end{equation}
and, after every such fixed derivative $\mathcal D$,
\begin{equation}
 \sup_{\mathbf z\in\mathcal T_L}
 |\mathcal D\mathcal E_K(\mathbf z)|
 \ll_{A,C_K,c_K}
 L^{C_K+C_0}T^{-2/\log_3T}\log_3T+L^{-A}
 \label{eq:parametric-HLZ-error}
\end{equation}
for every fixed $A>0$, where $C_0$ is an absolute fixed exponent coming from
the three zeta factors and the short contour.  In particular, any fixed
polylogarithmic loss in the auxiliary factor preserves the arbitrary
logarithmic saving in the HLZ diagonalization.
\end{lemma}

\begin{proof}
The Euler factorization in the proof of the diagonal lemma of
Heap--Li--Zhao~\cite{HLZ}, specialized to $h=k=1$ and with their second shift
replaced by $-z_2$, is exactly $D_{\mathbf z}(s)$. On the truncated rectangle,
the same zero-free region used there excludes zeros of the denominator
$\zeta(1-z_2+2s)$. Thus the only poles crossed when the $s$-line is moved left
are $s=0$ and the two shifted zeta poles
\[
 2s=z_2-z_1,\qquad 2s=z_2-z_3.
\]
The latter two are removable because of the assumed zeros of $K$, while the
residue at $s=0$ is $K_{T,\mathbf z}(0)D_{\mathbf z}(0)
=K_{T,\mathbf z}(0)Z_{\mathbf z,1,1}$.

We record the error estimate rather than invoking the HLZ argument as a black
box.  Put
\[
 V=B\sqrt{\log L},
\]
where $B$ is a fixed constant to be chosen after $A$ and $C_K$.  Truncate the
right vertical line at $|\Im s|=V$.  On that line the Dirichlet series is
absolutely convergent, while \eqref{eq:parametric-HLZ-K-bound} gives Gaussian
vertical decay; standard polynomial vertical bounds for the finitely many zeta
factors are therefore dominated by $e^{-c_Kv^2}$.  Choosing $B$ sufficiently
large makes the two discarded tails $O_A(L^{-A})$.

Move the truncated contour to $\Re s=-1/\log_3T$, exactly the narrow
zero-free-region displacement used in the proof of the HLZ diagonal lemma.
On the resulting rectangle their zero-free-region estimate gives, uniformly
for $\mathbf z\in\mathcal T_L$,
\[
 \zeta(1+z_1-z_2+2s),\quad
 \zeta(1-z_2+z_3+2s),\quad
 \zeta(1-z_2+2s)^{-1}
 \ll L^{C_0},
\]
with a fixed absolute exponent after the three factors are combined.  On the
left edge,
\[
 |T^{2s}|=T^{-2/\log_3T},
 \qquad
 |s|^{-1}\ll \log_3T
\]
near the only relevant small-height portion, while the Gaussian integral of
$(1+|v|)^{C_K}e^{-c_Kv^2}$ is bounded.  Hence the left edge is
\[
 \ll L^{C_K+C_0}T^{-2/\log_3T}\log_3T.
\]
On the two horizontal edges the factor $e^{-c_KV^2}$ beats every prescribed
power of $L$ once $B$ is enlarged, so their total contribution is
$O_A(L^{-A})$.  This proves \eqref{eq:parametric-HLZ-diagonal} and
\eqref{eq:parametric-HLZ-error}.

The same argument applies after any derivative $\mathcal D$ in the stated
finite family: these derivatives act only on $K$, the bound
\eqref{eq:parametric-HLZ-K-bound} is assumed for the differentiated auxiliary,
and the scalar Euler factor $D_{\mathbf z}$ is unchanged.  This yields the
uniform differentiated estimate.
\end{proof}

\begin{lemma}[Packet-weighted scalar HLZ diagonalization]
\label{lem:packet-HLZ-diagonal}
Let $\mathbf z\in\mathcal T_L$ and let $H_{\mathbf z}$ be the explicit
auxiliary function \eqref{eq:HLZ-explicit-auxiliary}.  Let
$W(u,\mathbf z)$ be a packet Mellin weight of
Lemma~\ref{lem:smooth-mellin-class}. Put
$\log_3T:=\log\log\log T$.  Pointwise for
$\mathbf z\in\mathcal T_L$, the invariant-source specialization $h=k=1$ of the
HLZ diagonal contour with auxiliary factor $H_{\mathbf z}W$ satisfies
\begin{equation}
\boxed{
W(0,\mathbf z)Z_{\mathbf z,1,1}
=\widetilde Z_{\mathbf z,1,1}[H_{\mathbf z}W]
+\mathcal E_{\rm diag}(\mathbf z;1,1),
}
\label{eq:packet-HLZ-diagonal}
\end{equation}
where, after any fixed number of the packet derivatives
$\partial_\rho,\partial_p$ used below,
\begin{equation}
\sup_{\mathbf z\in\mathcal T_L}
|\partial_\rho^r\partial_p^j\mathcal E_{\rm diag}(\mathbf z;1,1)|
\ll L^{C_W}
\left[
T^{-2/\log_3T}
(\log T)^C\log_3T+(\log T)^{-A}
\right]
\label{eq:packet-HLZ-error}
\end{equation}
for every fixed pair $r,j$ required later and every prescribed fixed $A>0$.  The physical shift derivatives are taken
only afterwards by the outer Cauchy integrals of
Proposition~\ref{prop:HLZ-three-face-Cauchy}; no derivative of
$H_{\mathbf z}$ across the collision locus is used.
In particular, for every fixed $A,D>0$,
\begin{equation}
\boxed{
T^{-2/\log_3T}L^D\log_3T\ll_{A,D}L^{-A}.
}
\label{eq:HLZ-superlog-saving}
\end{equation}
Thus the first term in \eqref{eq:packet-HLZ-error} is smaller than every fixed
negative power of $L$; after requesting the arbitrary exponent in the second
term large enough to absorb the fixed Cauchy and packet losses, the complete
error has the required arbitrary logarithmic saving.
\end{lemma}

\begin{proof}
The packet weight has the form
\[
W(u,\mathbf z)=\int_0^{C_*}a_{\rho,p,\mathbf z,T}(t)e^{Lut}\,dt
\]
(or a finite product of such factors). Lemma~\ref{lem:smooth-mellin-class}
shows that every fixed $\partial_\rho,\partial_p$ derivative has the same
polylogarithmic $L^1$ bounds. On every strip
\[
-\frac1{\log_3T}\le\Re u\le\frac{c_+}{L}
\]
with fixed $c_+>0$, these weights are bounded by a fixed power of $L$, because
$|e^{Lut}|\le e^{c_+C_*}$ on the right and is at most $1$ on the left.  By
\eqref{eq:HLZ-shift-separation}, the denominator in
\eqref{eq:HLZ-explicit-auxiliary} is bounded below by a fixed multiple of
$L^{-2}$.  Hence $H_{\mathbf z}(u)$ is bounded there by a fixed polynomial in
$L(1+|u|)$ times
\[
|e^{u^2}|=e^{(\Re u)^2-(\Im u)^2}.
\]
This gives exactly the Gaussian vertical decay required in the proof of the
HLZ diagonal lemma, uniformly on $\mathcal T_L$.

For clarity, \eqref{eq:HLZ-superlog-saving} follows on taking logarithms:
\[
\log\!\left(T^{-2/\log_3T}L^{A+D}\log_3T\right)
=-\frac{2\log T}{\log_3T}+(A+D)\log L+\log\log_3T\longrightarrow-\infty.
\]
Here $\log T/\log_3T\gg_{A,D}\log L$, so this estimate remains true after all
fixed polylogarithmic losses in \eqref{eq:packet-HLZ-error} are included.

The two possible shifted zeta poles are at
$2u=z_2-z_1$ and $2u=z_2-z_3$ and are cancelled pointwise by
\eqref{eq:HLZ-auxiliary-properties}.  Thus
\[
F_{T,\mathbf z}(u):=H_{\mathbf z}(u)W(u,\mathbf z)
\]
satisfies the hypotheses of Lemma~\ref{lem:parametric-HLZ-diagonal}, with a
fixed exponent $C_K$ large enough to absorb the packet and auxiliary polynomial
losses. Since $H_{\mathbf z}(0)=1$, its main residue is
\[
F_{T,\mathbf z}(0)Z_{\mathbf z,1,1}
=W(0,\mathbf z)Z_{\mathbf z,1,1}.
\]
Equation \eqref{eq:parametric-HLZ-error}, together with
\eqref{eq:HLZ-superlog-saving}, then gives
\eqref{eq:packet-HLZ-diagonal}--\eqref{eq:packet-HLZ-error}.

Finally, the required physical shift coefficients are finite Cauchy residues
of the three-face identity \eqref{eq:HLZ-three-face-Cauchy}.  On the separated
tori the Cauchy denominators cost only fixed powers of $L$, while the HLZ
arbitrary-log saving may be requested with a correspondingly larger exponent
before the finite residues are taken.  Thus the same displayed error class
survives the finite shift extraction without ever evaluating
$H_{\mathbf z}$ at a colliding triple.
\end{proof}

\begin{lemma}[Local-height rescaling of the HLZ diagonal]
\label{lem:HLZ-local-height-rescaling}
Fix a retained physical height $\tau\asymp T$ and put
\[
X(\tau):=\frac{\tau}{2\pi}.
\]
In the scalar diagonal contour of Lemma~\ref{lem:packet-HLZ-diagonal}, replacing
its fixed factor
\[
T^{2u}(mn)^{-u}
\]
by
\[
X(\tau)^{2u}(mn)^{-u}
\]
does not change the $u=0$ main residue and preserves the error class
\eqref{eq:packet-HLZ-error}, uniformly on every retained height block.  The same
is true after the fixed packet, curvature, and physical-shift Cauchy
derivatives used below.  Moreover, if the logarithmic endpoint $U$ is tied to
physical height by $\tau=2\pi e^{LU}$, then the first two total endpoint
derivatives
\[
D_U:=\partial_U+L\tau\partial_\tau
\]
also preserve the same arbitrary-log error class.
\end{lemma}

\begin{proof}
Write
\[
R_\tau(u):=\left(\frac{X(\tau)}{T}\right)^{2u}.
\]
Algebraically, the diagonal contour with the local scale $X(\tau)$ and packet
weight $W$ is the original $T$-scale contour with packet weight $R_\tau W$.
Since $\tau\in[T,2T]$, the ratio $X(\tau)/T$ stays in a fixed compact subset of
$(0,\infty)$.  Hence on the whole strip
\[
-\frac1{\log_3T}\le \Re u\le\frac{c_+}{L}
\]
one has
\[
|R_\tau(u)|\ll1,
\qquad
(\tau\partial_\tau)^jR_\tau(u)=(2u)^jR_\tau(u),
\qquad
(L\tau\partial_\tau)^jR_\tau(u)=(2Lu)^jR_\tau(u)
\]
for every fixed $j$.  Thus multiplication by $R_\tau$, even after the first
two $L\tau\partial_\tau$ derivatives, preserves the Gaussian vertical decay
and merely enlarges the fixed polylogarithmic exponent in
Lemma~\ref{lem:parametric-HLZ-diagonal}.

Most importantly,
\[
R_\tau(0)=1.
\]
Therefore the residue produced when the $u$-contour crosses the origin is still
\[
(R_\tau H_{\mathbf z}W)(0,\mathbf z)Z_{\mathbf z,1,1}
=W(0,\mathbf z)Z_{\mathbf z,1,1}.
\]
For every $j\ge1$,
\[
(L\tau\partial_\tau)^jR_\tau(0)=0,
\]
so differentiating the scale multiplier contributes no new main residue at the
origin.  On the displaced contour the factors $(Lu)^j$ are absorbed by the
Gaussian vertical decay and by increasing the fixed polylogarithmic exponent
in Lemma~\ref{lem:parametric-HLZ-diagonal}.  The $\partial_U$ part of $D_U$ acts
only on the pre-existing packet/endpoint weight and is already covered by the
packet derivative hypotheses.  Hence the first two total endpoint derivatives,
the left and horizontal contour estimates, and then the finite Cauchy residues
retain the same arbitrary-log error class \eqref{eq:packet-HLZ-error}.  This
proves the claimed local-height rescaling.
\end{proof}

\begin{lemma}[Auxiliary-shift rigidity of the finite-model packet map]
\label{lem:aux-shift-rigidity-packet-map}
Fix a retained height block and restrict to the principal stationary part of the
finite model source $H_{\rm mod}$.  Before the HLZ $u$-contour deformation, the
three auxiliary Cauchy variables $\mathbf z=(z_1,z_2,z_3)$ act only on a
packet-label-independent scalar amplitude.  More precisely, for every fixed
retained physical shift $p$ and every basis packet $e_r$ there are a one-sided
interval $I_r(p)$ and a function $b_{r,T,p}$, both independent of $\mathbf z$,
such that every packet-pair coefficient of the completed finite-model
three-face source has the form
\begin{equation}
\boxed{
\mathbf 1_{I_r(p)\cap I_s(p)}(t)
 b_{s,T,p}(t)\overline{b_{r,T,p}(t)}\,
 c_{T,p}(t;u,\mathbf z),
}
\label{eq:preHLZ-shift-rigid-factorization}
\end{equation}
where all dependence on $u,\mathbf z$ is contained in the common scalar factor
$c_{T,p}$.  Consequently, for each fixed $p$, the principal packet analysis map
\[
(\mathscr C_{T,p}e_r)(t):=\mathbf 1_{I_r(p)}(t)b_{r,T,p}(t)
\]
is independent of the auxiliary Cauchy shifts; in particular, in this exact
pre-HLZ representation,
\begin{equation}
\boxed{\partial_{z_j}\mathscr C_{T,p}=0\qquad(j=1,2,3).}
\label{eq:C-map-shift-rigidity}
\end{equation}
The statement concerns the principal finite-model source only.  Higher
stationary corrections and nonstationary leakage remain in the regular ledger
of Proposition~\ref{prop:regular-ledger}.
\end{lemma}

\begin{proof}
The chronology of the exact source fixes the dependence.  In the invariant
contour form the packet factor is already
\[
\mathcal B_{e_r,e_s}(s)=\psi_s(z(s))\psi_r^\#(z(s))
\]
by \eqref{eq:packet-polarization}; it contains neither the physical Hardy shift
nor any HLZ auxiliary Cauchy variable.  The finite replacement
$H_{\rm mod}=F_b-2P+P(s-2/L)$ is a scalar logarithmic-derivative source, so
multiplying by it cannot alter the packet labels.  The variables
$z_1,z_2,z_3$ are introduced only afterwards by the exact scalar three-face
Cauchy identity \eqref{eq:HLZ-three-face-Cauchy}.  Thus, before any stationary
reduction, every $\mathbf z$-dependent shifted-zeta, Cauchy-denominator and
archimedean factor is common to all packet labels.

After Dirichlet expansion, a small auxiliary zeta shift changes a Dirichlet
coefficient by factors such as $n^{-z_j}$ and changes the functional-equation
factor by a ratio of shifted gamma factors.  We keep all such factors in the
scalar amplitude.  Equivalently, a shifted one-$\chi$ kernel may be written as
\[
A_{\mathbf z}(t,\xi)\,
\chi(1-c-it)\xi^{-c-it},
\]
where $A_{\mathbf z}$ is packet-label independent and, on the separated
$|z_j|\asymp L^{-1}$ tori, has only fixed polylogarithmic derivative loss.  The
rapid phase is therefore exactly the unshifted one-$\chi$ phase of
Section~\ref{sec:stationary}.  In particular its saddle, saddle cutoff and
principal Hessian normalization are independent of $\mathbf z$.

The only transformation of the packet factor at principal stationary level is
therefore the exact one-sided dilation
\[
\mathcal S_w^{\rm pr}=D_w\mathcal S_0^{\rm pr}
\]
of \eqref{eq:dilation-covariance}.  Here $w$ is the logarithmic packet translation, fixed by the packet coordinate,
physical carrier and the already-fixed physical shift $p$, not by the auxiliary
Cauchy shifts.  Hence the transported packet support $I_r(p)$ and the one-sided
factor $b_{r,T,p}$ are independent of $\mathbf z$.  All shifted-zeta, gamma,
Cauchy and Mellin factors remain in one common scalar amplitude
$c_{T,p}(t;u,\mathbf z)$.  This proves
\eqref{eq:preHLZ-shift-rigid-factorization}, and
\eqref{eq:C-map-shift-rigidity} follows directly from the definition of
$\mathscr C_{T,p}$.
\end{proof}

\begin{lemma}[Common-output lift of the scalar HLZ contour error]
\label{lem:HLZ-error-common-output}
Fix a retained physical shift $p$ and a height block.  Suppose a principal
finite-model packet weight has the pre-HLZ factorization of
Lemma~\ref{lem:aux-shift-rigidity-packet-map}, so that
\begin{equation}
W_{rs,L}(u,\mathbf z)
=\int
 (\mathscr C_{T,p}e_s)(t)
 \overline{(\mathscr C_{T,p}e_r)(t)}
 F_{T,p,t}(u,\mathbf z)\,dt,
\label{eq:HLZ-error-slice-factorization}
\end{equation}
with $\mathscr C_{T,p}$ independent of $\mathbf z$.  Assume uniformly in $t$
that $H_{\mathbf z}F_{T,p,t}$ satisfies the hypotheses of
Lemma~\ref{lem:parametric-HLZ-diagonal}.  Then the scalar HLZ contour error is
not merely entrywise small: there is a packet-label-independent multiplication
operator $M_{\varepsilon,p,\mathbf z}$ on the same output space such that
\begin{equation}
\boxed{
\mathcal E_{rs}(p,\mathbf z)
=\left\langle
 M_{\varepsilon,p,\mathbf z}\mathscr C_{T,p}e_s,
 \mathscr C_{T,p}e_r
\right\rangle.
}
\label{eq:HLZ-error-common-output}
\end{equation}
For every prescribed fixed $A>0$, after increasing the scalar saving exponent
if necessary,
\begin{equation}
\boxed{
\|M_{\varepsilon,p,\mathbf z}\|_{\rm op}
\ll_A L^{-A+C}
}
\label{eq:HLZ-error-multiplier-bound}
\end{equation}
uniformly on the separated Cauchy torus.  After any fixed finite Cauchy residue
in $\mathbf z$, the resulting error is still a pullback through the same
$\mathscr C_{T,p}$, with only a fixed additional power of $L$ in the multiplier
bound.  No packet-dimension factor is lost.
\end{lemma}

\begin{proof}
For fixed $\mathbf z$, let $\mathfrak E_{T,\mathbf z}[F]$ denote the error
functional in the proof of Lemma~\ref{lem:parametric-HLZ-diagonal}: the sum of
the displaced left-edge integral, the two horizontal connectors and the
Gaussian truncation tails.  Each term is linear in the auxiliary analytic
weight $F$.  Moreover the Gaussian truncation makes the relevant integrals
absolutely convergent, uniformly for the slice family in
\eqref{eq:HLZ-error-slice-factorization}.  Hence Fubini gives
\[
\mathcal E_{rs}(p,\mathbf z)
=\int
 (\mathscr C_{T,p}e_s)(t)
 \overline{(\mathscr C_{T,p}e_r)(t)}
 \mathfrak E_{T,\mathbf z}[H_{\mathbf z}F_{T,p,t}]\,dt.
\]
Define
\[
\varepsilon_{p,\mathbf z}(t)
:=\mathfrak E_{T,\mathbf z}[H_{\mathbf z}F_{T,p,t}].
\]
This proves \eqref{eq:HLZ-error-common-output} with
$M_{\varepsilon,p,\mathbf z}$ multiplication by
$\varepsilon_{p,\mathbf z}(t)$.

The uniform slice hypotheses and Lemma~\ref{lem:parametric-HLZ-diagonal} give
\[
\sup_t|\varepsilon_{p,\mathbf z}(t)|
\ll
L^C T^{-2/\log_3T}\log_3T+L^{-A'}.
\]
Choose $A'$ larger than $A$ by enough to absorb the fixed slice and auxiliary
losses, and use \eqref{eq:HLZ-superlog-saving}.  This proves
\eqref{eq:HLZ-error-multiplier-bound}.  Finally, by
Lemma~\ref{lem:aux-shift-rigidity-packet-map} the map $\mathscr C_{T,p}$ is
independent of $\mathbf z$.  Therefore a finite Cauchy residue differentiates
or integrates only the scalar multiplier.  The separated torus contributes
only fixed powers of $L$, so the same common-output representation survives the
residue extraction.  Since the estimate is on the multiplication-operator norm
before pullback, no factor depending on $|\mathcal I_T|$ appears.
\end{proof}

\begin{lemma}[Exact pre-HLZ packet separation of the finite-model principal source]
\label{lem:pre-HLZ-packet-separation}
Fix one retained height block and one right-edge orientation.  Remove the
$f'$-class by \eqref{eq:exact-fprime-separation}, freeze the scalar half-self
factor at $F_b$, and retain the finite-model principal stationary term.  Before
any HLZ $u$-contour deformation and before the finite physical-shift Cauchy
residues, the packet dependence is carried by one shift-independent analysis
map.

More precisely, for basis packets $e_r,e_s$ there are functions
\[
a_{s,T}(x):=\phi_L(x)e^{-i\tau_sx}
\]
and a packet-label-independent scalar lag kernel
$\kappa_T(d;u,\mathbf z,p)$ such that the pre-HLZ principal slice is exactly
\begin{equation}
\boxed{
\mathcal Q_{rs}^{\rm preHLZ}(u,\mathbf z,p)
 =\iint a_{s,T}(x)\overline{a_{r,T}(y)}
 \kappa_T(y-x;u,\mathbf z,p)\,dx\,dy .
}
\label{eq:pre-HLZ-lag-factorization}
\end{equation}
All dependence on the HLZ Cauchy shifts $\mathbf z$, the curvature parameter
$p$, the diagonal Mellin variable $u$, the frozen block data, and the scalar
Dirichlet coefficients lies in $\kappa_T$; none of these parameters occurs in
$a_{r,T}$.

Consequently, after Fourier diagonalization of the lag operator, there are one
Hilbert output space $\mathcal H_{\rm fib}$, one packet analysis map
\[
\mathscr C_T:\mathbb C^{\mathcal I_T}\to\mathcal H_{\rm fib},
\]
independent of $(u,\mathbf z,p)$, and packet-label-independent multiplication
operators $M_{u,\mathbf z,p}$ such that
\begin{equation}
\boxed{
\mathcal Q_{rs}^{\rm preHLZ}(u,\mathbf z,p)
 =\left\langle M_{u,\mathbf z,p}\mathscr C_Te_s,
 \mathscr C_Te_r\right\rangle .
}
\label{eq:pre-HLZ-common-output}
\end{equation}
Every fixed $p$- or auxiliary-shift derivative acts only on
$M_{u,\mathbf z,p}$.  The flat-core restriction is a parameter-independent
restriction of the two packet legs; the complementary one-leg and two-leg
reserve/cap terms are exactly the exterior-support terms of
Proposition~\ref{prop:exterior-packet-support} and are not part of the frozen
principal source.
\end{lemma}

\begin{proof}
Put $s=c+it$ and $\sigma_0=c-1/2$.  From the packet definition
\eqref{eq:packet-polarization} and the Fourier representation
\[
\psi_{L,\tau}(z)=\int\phi_L(x)e^{ix(z-\tau)}\,dx,
\]
using that $\phi_L$ is real, one obtains the exact basis-packet identity
\begin{equation}
\mathcal B_{e_r,e_s}(c+it)
 =\iint\phi_L(x)\phi_L(y)e^{-i\tau_sx+i\tau_ry}
 e^{\sigma_0(x-y)}e^{it(x-y)}\,dx\,dy .
\label{eq:packet-right-edge-double-Fourier}
\end{equation}
In particular the packet centres occur only in the two separated phases.
Neither the normalized curvature shift $p$ nor the later HLZ Cauchy shifts
occur in \eqref{eq:packet-right-edge-double-Fourier}, because the invariant
polarization is evaluated at the unshifted central point $s$.

Now take one scalar Dirichlet monomial $\xi^{-c-it}$ in the frozen finite-model
source.  If $d=y-x$, then
\begin{equation}
 e^{\sigma_0(x-y)}e^{it(x-y)}\xi^{-c-it}
 =e^{d/2}(\xi e^d)^{-c-it}.
\label{eq:packet-lag-one-chi}
\end{equation}
Thus, apart from the two packet phases, the entire one-$\chi$ oscillatory
integrand depends on $(x,y)$ only through the difference $d$.  The principal
stationary calculation of Section~\ref{sec:stationary} has
\[
\eta=\xi e^d,
\qquad t_*=2\pi\eta,
\]
and by \eqref{eq:stationary-dilation} the power of $\eta$ in the Stirling
amplitude cancels exactly against the Gaussian Hessian factor.  Hence no
factor depending on $x+y$, or on either packet leg separately, is generated.
The smooth height cutoff is transported by the same exact dilation
\eqref{eq:dilation-covariance}.  Every retained archimedean, curvature,
Cauchy-shift, translated-face, and frozen block factor is scalar; when evaluated
at the saddle it is therefore a function of $d$ and the displayed scalar
parameters only.  The summation over scalar Dirichlet monomials is legitimate
on the original contour without analytic continuation: by
\eqref{eq:symmetric-contour-rectangle} the right edge is $c=2$, and for all
large $T$ the translated face still has
\[
\Re(s-2/L)=2-2/L>1.
\]
Hence both $P(s)$ and $P(s-2/L)$, and the finite products used in the frozen
finite model, are absolutely convergent there, uniformly on the compact
$d$-range generated by $\supp\phi_L$.  Their fixed $p$- and auxiliary-shift
derivatives cost only polylogarithmic factors; the uniform stationary-symbol
bounds of Lemma~\ref{lem:stationary-symbol} apply after the saddle evaluation.
Thus the summed kernel is a bounded smooth lag kernel on that compact range and
\eqref{eq:pre-HLZ-lag-factorization} follows.

Insert a fixed smooth cutoff equal to one on the compact difference set
generated by $\supp\phi_L$.  The resulting lag kernel therefore defines a
bounded convolution operator $T_{u,\mathbf z,p}$ on the packet Fourier coordinate.
Fourier transform diagonalizes it:
\[
T_{u,\mathbf z,p}=\mathcal F^{-1}M_{u,\mathbf z,p}\mathcal F.
\]
Taking $\mathscr C_Te_s=\mathcal F a_{s,T}$, together with the common normalized
one-sided stationary factors already present in the principal source, gives
\eqref{eq:pre-HLZ-common-output}.  In the diagonal Mellin/Perron representation
this multiplier is precisely the common-prefix scalar weight; its $u$-dependence
is the factor $e^{Lut}$ appearing in \eqref{eq:model-Fubini-weight} below.
Therefore differentiation in $p$ or in the auxiliary Cauchy variables acts on
the output multiplier and never on $\mathscr C_T$.

Finally split each packet leg by the fixed central flat-core restriction and
its complement.  The central--central product is the frozen principal source;
the central--exterior, exterior--central, and exterior--exterior products are
exactly the one-leg and two-leg terms assigned by
Proposition~\ref{prop:exterior-packet-support}.  Since this split is made by
parameter-independent restrictions before differentiation, it creates no
endpoint or packet-analysis jets.  This proves the lemma.
\end{proof}

\begin{lemma}[Fubini lift of the finite-model three-$Z$ faces]
\label{lem:model-HLZ-Fubini-lift}
For the logarithmic-derivative faces $-2P+P(s-2/L)$ of the finite model source
$H_{\rm mod}$ in \eqref{eq:Hmod}, retain the packet polarization before the
scalar HLZ diagonalization.  Fix a retained height block and a packet pair, let
$I_r,I_s$ be their two one-sided logarithmic support intervals after the
principal stationary dilation, and put $J_{rs}:=I_r\cap I_s$.  Before the
scalar deformation the packet-pair weight has the one-sided form
\begin{equation}
\boxed{
W_{rs,L}(u,\mathbf z)
=\int_{J_{rs}} b_{s,T}(t)\overline{b_{r,T}(t)}
 c_T(t,\mathbf z)e^{Lut}\,dt,
}
\label{eq:model-Fubini-weight}
\end{equation}
where $b_{r,T}$ is independent of the three auxiliary Cauchy shifts and all
shift dependence is contained in the packet-label-independent scalar factor
$c_T$. Consequently there is a shift-independent packet analysis map
$\mathscr C_T$ such that the scalar three-$Z$ model main term and its scalar HLZ
contour error, after the finite physical-shift Cauchy residues, are both
pullbacks through $\mathscr C_T$ of packet-label-independent multiplication
operators.
Moreover, for every prescribed fixed $A>0$, the error multiplier and the fixed
endpoint/curvature derivatives used below are $O(L^{-A+C})$ for a fixed $C$.
\end{lemma}

\begin{proof}
Lemma~\ref{lem:pre-HLZ-packet-separation} gives the exact factorization before
the scalar HLZ contour is moved: the packet labels are carried by one
$(u,\mathbf z,p)$-independent analysis map, while the shifted zeta,
archimedean, curvature, Cauchy, and diagonal Mellin factors act only on the
common output fibre.  Writing that multiplier in the diagonal Mellin/Perron
coordinate is exactly \eqref{eq:model-Fubini-weight}.

For fixed fibre coordinate $t$, the remaining analytic weight is a scalar
multiple of $e^{Lut}$. On the strip
$-1/\log_3T\le\Re u\le c_+/L$ it is uniformly bounded on the right and gains
exponential decay on the left; the fixed derivatives of the remaining smooth
scalar factors cost only powers of $L$. Hence the proof of
Lemma~\ref{lem:packet-HLZ-diagonal} applies uniformly to each such scalar
slice. The HLZ contour identity is linear in its auxiliary analytic function,
so Fubini gives the packet matrix as
$\mathscr C_T^*M\mathscr C_T$, both for the scalar main term and for the scalar
contour error. The separated finite Cauchy residues still act only on the
common scalar multiplier and cost a fixed power of $L$. Since the scalar HLZ
saving exponent is arbitrary, increasing it absorbs every fixed derivative and
Cauchy loss and proves the stated error bound.
\end{proof}

\begin{hypothesis}[Uniform-twisted packet local-projection bridge]
\label{hyp:packet-bilinear-HLZ}
For the invariant packet source, the scalar identities
\eqref{eq:HLZ-three-face-Cauchy} and
\eqref{eq:packet-HLZ-diagonal} admit the following carrier-preserving lift for
the exact right-edge source \eqref{eq:exact-fprime-separation} after the
$f'$-class has been removed.

For every retained height block, every $0\le q\le2$, and every exact HLP level
$0\le k\le K_0$, first form the Dirichlet series on the fixed
absolute-convergence line and take the required curvature derivative at
$p=0$.  Then:
\begin{enumerate}[label=(\roman*)]
\item the resulting completed coefficient is exactly the finite Cauchy transform
of the three-face HLZ master in Proposition~\ref{prop:HLZ-three-face-Cauchy},
with the packet Mellin weight inserted before the $u$-residue; in particular
there is no unrecorded face, sign, normalization, or source coefficient between
\eqref{eq:exact-fprime-separation} and the completed coefficient
$\mathcal F_{k,L}^{[q]}$ used in Section~\ref{sec:hlp-local};
\item there are a Hilbert output space $\mathcal H_{\rm out}$ and one bounded
analysis map
\[
\mathscr C_T:\mathbb C^{\mathcal I_T}\longrightarrow\mathcal H_{\rm out},
\]
independent of the three Cauchy shifts and of the packet pair, such that every
completed low-level quadratic form before the $u$-residue is
$\mathscr C_T^*M_{k,q,T}(u,\mathbf z)\mathscr C_T$, with all shift dependence
in the output operator $M_{k,q,T}$; the total internal HLP Dirichlet index is
summed inside this output operator and does not create an additional external
packet fibre;
\item after the pointwise scalar HLZ deformation on the separated torus and the
levelwise Perron deformation, every common contour-error term remains on the
same output channel and, on each retained unit shell, has a finite sum of
operator-valued kernels of the form
\begin{equation}
D_{j,T}^{\rm err}(t,s;p)
=\frac1L\,a_{j,T}(t,s;p)\,
 h_{T}^{\rm err}\!\left(\nu_T-\frac{|\log(J_j+t)-\log(J_j+s)|}{L},p\right).
\label{eq:hyp-common-error-kernel}
\end{equation}
The amplitude acts on a packet-independent local source/refinement fibre
$\mathcal K_j$ with $\dim\mathcal K_j\le L^{C_{\rm src}}$ and, through the
fixed derivative orders used below,
\[
\|\partial_t^a\partial_s^b\partial_p^q a_{j,T}(t,s;p)\|_{\op}
\ll L^{C_{\rm src}}J_j^{-a-b}
\qquad(a+b\le2,\ q\le2).
\]
For every prescribed fixed $A>0$, the scalar error profile may be taken so that
\[
\max_{0\le r,q\le2}
|\partial_\rho^r\partial_p^q h_T^{\rm err}(\rho,p)|
\ll L^{-A+C_A}.
\]
The number of summands and all constants are bounded by fixed powers of $L$ and
are independent of $\#\mathcal I_T$.
\end{enumerate}
Equivalently, the HLZ/Perron transfer is assumed at the quadratic-form/output
level before scalar contraction, not entrywise after a packet matrix has been
formed.
\end{hypothesis}

\begin{remark}[Optional stronger exact-HLP route]
Hypothesis~\ref{hyp:packet-bilinear-HLZ} is stronger than the scalar theorem of
Heap--Li--Zhao and is retained only to document the earlier levelwise
exact-HLP completion route.  The proof of Theorem~\ref{thm:main} does not use
this hypothesis: the finite model source is lifted by
Lemma~\ref{lem:model-HLZ-Fubini-lift}, while the exact--model difference is kept
as a direct one-$\chi$ family on the absolute-convergence right edge.  Any
lemma below that explicitly assumes Hypothesis~\ref{hyp:packet-bilinear-HLZ}
belongs only to that optional comparison route.
\end{remark}

\begin{lemma}[Normalization of the finite HLZ Euler factor]
\label{lem:HLZ-finite-factor}
In the notation of the completed HLZ diagonal, the finite local factor attached
to $k$ is
\begin{equation}
\mathcal E_{\alpha,\beta,\gamma}(k)
:=\prod_{p^{\kappa_p}\parallel k}
\frac{\displaystyle\sum_{m\ge0}
 f_{\alpha,\gamma}(p^{m+\kappa_p})p^{-m(1+\beta)}}
{\displaystyle\sum_{m\ge0}
 f_{\alpha,\gamma}(p^m)p^{-m(1+\beta)}},
\label{eq:HLZ-finite-factor}
\end{equation}
where
\[
f_{\alpha,\gamma}=\mu* n^{-\alpha}*n^{-\gamma}
\]
in Dirichlet-convolution notation. Then
\begin{equation}
\boxed{\mathcal E_{0,0,0}(k)=1.}
\label{eq:HLZ-factor-zero}
\end{equation}
Moreover, uniformly for $k\le Y=L^\kappa$ and
$|\alpha|+|\beta|+|\gamma|\le C/L$,
\begin{equation}
\boxed{
\mathcal E_{\alpha,\beta,\gamma}(k)
=1+O_C\!\left(\frac{\log Y}{L}\right),
}
\label{eq:HLZ-factor-small}
\end{equation}
and every fixed derivative with respect to the normalized variables
$L\alpha,L\beta,L\gamma$ is bounded by a fixed power of
$\log Y/L$ (with the zeroth derivative bounded uniformly).
\end{lemma}

\begin{proof}
For one prime put $z=p^{-(1+\beta)}$. The local generating series is
\[
\sum_{m\ge0}f_{\alpha,\gamma}(p^m)z^m
=\frac{1-z}{(1-p^{-\alpha}z)(1-p^{-\gamma}z)}.
\]
At $\alpha=\beta=\gamma=0$ this is $(1-z)^{-1}$, so every local coefficient is
$1$ and the numerator and denominator in
\eqref{eq:HLZ-finite-factor} coincide. This proves
\eqref{eq:HLZ-factor-zero}.

On the stated shift polydisc one has $|z|\le p^{-1+o(1)}$, uniformly for
$p\le Y$. Hence the local series and its reciprocal are analytic and uniformly
bounded on a fixed slightly larger polydisc. Differentiating a local factor
$j$ times in the unnormalized shifts costs at most
$O_j((\kappa_p+1)^j(\log p)^j)$. Summing logarithmic derivatives over
$p^{\kappa_p}\parallel k$ and using
\[
\sum_{p^{\kappa_p}\parallel k}\kappa_p\log p=\log k\le\log Y
\]
shows that the first logarithmic derivative is $O(\log Y)$ and every fixed
higher derivative is bounded by a fixed power of $\log Y$. Passing to the
normalized variables divides each derivative by $L$. The mean-value theorem,
together with \eqref{eq:HLZ-factor-zero}, gives
\eqref{eq:HLZ-factor-small} and the asserted derivative bounds.
\end{proof}

\begin{proposition}[Height versus fibre endpoint]
\label{prop:X-tagging}
In the completed three-$Z$ source, every explicit power of the interpolation
height $X$ belongs to the functional-equation/archimedean factor, whereas the
packet-dependent finite endpoint is produced by the diagonal Mellin factor.
More precisely,
\[
X_{\rm arch}=\frac{t}{2\pi},
\qquad
\lambda(t):=\frac{\log X_{\rm arch}}L=1+O(L^{-1})
\]
on a retained dyadic height block.  The original HLZ diagonal is written with
the fixed scale $T$ and hence contains
\[
T^{2u}(mn)^{-u}
=\left(\frac{T^2}{hk\ell^2}\right)^u.
\]
Lemma~\ref{lem:HLZ-local-height-rescaling} replaces this, at the level of the
same $u=0$ principal residue, by the local-scale factor
\[
X_{\rm arch}^{2u}(mn)^{-u}
=\left(\frac{X_{\rm arch}^2}{hk\ell^2}\right)^u,
\]
with only the already-recorded HLZ contour error.  Thus the diagonal Mellin
factor produces a moving fibre endpoint at the same local height.  After the
two Hermitian stationary cutoffs are intersected, the full-fibre
scale is
\[
Y_{\rm fib}:=X_{\rm arch}=\frac{t}{2\pi},
\qquad
\frac{\log Y_{\rm fib}}L=\lambda(t)=1+O(L^{-1}).
\]
Then
\[
\ell\le Y_{\rm fib}e^{-|v|},
\qquad
v=w+\log(h/k),
\]
so
\begin{equation}
\boxed{
\rho(v)=\frac{\log Y_{\rm fib}-|v|}{L}
=1-\frac{|v|}{L}+O(L^{-1}).
}
\label{eq:rho-v}
\end{equation}
The symbol $Y_{\rm fib}\asymp T$ is unrelated to the arithmetic truncation
$Y=L^\kappa$ introduced later. This is the \emph{raw} stationary coordinate:
the term $\log(h/k)=\log h-\log k$ is removed later by an exact pair of packet
translations in the common positive analysis map, not by a small endpoint
approximation. Consequently packet lag changes the finite endpoint $\rho$ but
does not replace the explicit archimedean powers by powers of $\rho$.
\end{proposition}

\begin{proof}
Before the arithmetic master is contracted, the three-face interpolation
contains the common monomial
\[
X^{-(\alpha+\beta+\gamma)/2}
X^{(z_1-z_2+z_3)/2}.
\]
Taking the finite Cauchy residues gives the explicit factors
\[
1,\qquad X^{-\alpha-\beta},\qquad X^{-\beta-\gamma}.
\]
These factors are present even after the smooth Mellin cutoff has been
suppressed from the displayed fixed-fibre coefficient, so their $X$ is the
physical height $X_{\rm arch}=t/(2\pi)$.  After
\eqref{eq:hlz-diagonal} the original Mellin variable occurs through
$T^{2u}(hk\ell^2)^{-u}$.  Apply
Lemma~\ref{lem:HLZ-local-height-rescaling}: multiplying the admissible packet
weight by $(X_{\rm arch}/T)^{2u}$ changes neither its value at $u=0$ nor the
HLZ error class, and converts the principal diagonal factor exactly to
\[
\left(X_{\rm arch}^2/(hk\ell^2)\right)^u.
\]
Thus the local diagonal cutoff and the explicit archimedean powers are tagged
by the same physical height for two different reasons: the latter are present
in the Cauchy master itself, while the former is obtained from the HLZ diagonal
by the residue-preserving rescaling lemma.  The packet dilation of
Section~\ref{sec:stationary} changes the two oriented cutoffs to
$Y_{\rm fib}e^{-v}$ and $Y_{\rm fib}e^v$; their common part is
$Y_{\rm fib}e^{-|v|}$, proving \eqref{eq:rho-v}.

The distinction survives differentiation. For instance,
\[
\partial_\gamma
\left[X^{-q-\gamma}A^\gamma\right]_{\gamma=0}
=X^{-q}\log(A/X),
\]
so the $X^{-q}$ and the $-\log X$ term have the same archimedean provenance,
whereas the $A=h\ell$ contribution is a fibre variable. Likewise, for
$\delta=2/L$,
\[
\left(\frac{kX}{h n_2n_3}\right)^\delta
=\left(\frac kh\right)^\delta X^\delta(n_2n_3)^{-\delta},
\]
with
\[
X^\delta=e^{2\lambda}=e^2(1+O(L^{-1})),
\]
while only the last factor is integrated over the common fibre. Hence the
archimedean and endpoint roles are distinct already at the source level.
\end{proof}

\begin{lemma}[Exact outer-ratio translation intertwiner]
\label{lem:outer-translation-intertwiner}
On the raw stationary coordinate the common-prefix lag is
\[
v=w+\log h-\log k.
\]
On the ambient packet Fourier space $L^2(\mathbb R,du)$ let
\[
(U_hF)(u):=F(u-\log h),
\qquad
(U_kF)(u):=F(u-\log k).
\]
Then $U_h,U_k$ are unitary, and the raw pair kernel depending on
$w+\log(h/k)$ is exactly the pullback by $U_h^*U_k$ of the universal kernel in
$w$. This identity holds for the full smooth packet profile, including its
transition caps; no flat-amplitude approximation is involved. The reserve
condition \eqref{eq:packet-bandwidth-reserve} is used only later when the
physical stationary band is represented inside the flat plateau for shell
analysis.
\end{lemma}

\begin{proof}
Multiplication of a Dirichlet monomial by $h^{-it}$ translates the packet
Fourier coordinate by $\log h$; the adjoint $k$-orientation gives the analogous
translation by $\log k$. Translation on $L^2(\mathbb R)$ preserves Lebesgue
measure and the full translated packet profile, so the relative displacement is
exactly $\log h-\log k$ and conjugation by $U_h,U_k$ replaces
$w+\log(h/k)$ by $w$ without a cutoff error. These ambient unitaries are part
of the common packet analysis map in both $A_T^G$ and $\Gfr_T$. When the later
shell representation is restricted to the physical stationary band,
$|\log h|,|\log k|\le\kappa\log L$ and the chosen reserve keep the translated
physical copies strictly inside the constant plateau.
\end{proof}

\begin{lemma}[Smoothed AFE weights and sharp-fibre transfer]
\label{lem:J-AFE-sharp-transfer}
Let $|\alpha|+|\beta|\le C/L$, let $\tau\in[T,2T]$, and put
\[
X:=\frac{\tau}{2\pi},
\qquad
U_\tau:=\frac{\log X}{L},
\qquad
A:=L\alpha,\quad B:=L\beta.
\]
Choose once and for all an even entire function $G_{\rm AFE}$ with
$G_{\rm AFE}(0)=1$ and Gaussian decay on vertical strips; for example
$G_{\rm AFE}(u)=e^{u^2}$.  Define the pure Mellin cutoff
\begin{equation}
V(x):=\frac1{2\pi i}\int_{(c)}G_{\rm AFE}(u)x^{-u}\,\frac{du}{u},
\qquad c>0.
\label{eq:pure-AFE-weight}
\end{equation}
Then
\begin{equation}
\boxed{V(x)+V(x^{-1})=1\qquad(x>0).}
\label{eq:AFE-complement}
\end{equation}
Set
\[
\gamma(s):=\pi^{-s/2}\Gamma(s/2),\qquad
\Lambda(s):=\gamma(s)\zeta(s),
\]
and
\[
s_1:=\frac12+\alpha+i\tau,\qquad
s_2:=\frac12+\beta-i\tau.
\]
Thus $\Lambda(s)=\Lambda(1-s)$ and $\Lambda$ has simple poles at $0$ and
$1$.  After the local conductor $X$ is extracted, define the two smoothed
weights by
\begin{equation}
V_{\alpha,\beta,\tau}^{\pm}(x)
=\frac1{2\pi i}\int_{(c)}G_{\rm AFE}(u)
 R_{\alpha,\beta,\tau}^{\pm}(u)x^{-u}\,\frac{du}{u},
\label{eq:exact-AFE-weights}
\end{equation}
where
\begin{equation}
R_{\alpha,\beta,\tau}^{\pm}(0)=1,
\qquad
R_{\alpha,\beta,\tau}^{\pm}(u)
=1+O_C\!\left(\frac{(1+|u|)^2}{T}\right)
\quad (|\Im u|\le T^{1/4})
\label{eq:AFE-gamma-ratio}
\end{equation}
uniformly on every fixed real strip used in the contour shift.  On the
complementary part $|\Im u|>T^{1/4}$ the Gaussian factor
$G_{\rm AFE}(u)$ makes the complete Mellin integrand $O_M(T^{-M})$ for every
fixed $M$, uniformly in the same shifts.

The full shifted-product identity has four additional polar terms.  More
precisely, if
\[
\widetilde s_1:=1-s_2=\frac12-\beta+i\tau,
\qquad
\widetilde s_2:=1-s_1=\frac12-\alpha-i\tau,
\]
then
\begin{equation}
\boxed{
\begin{aligned}
\zeta(s_1)\zeta(s_2)={}&
\sum_{m,n\ge1}\frac{
 V^+_{\alpha,\beta,\tau}(mn/X)}{m^{s_1}n^{s_2}}\\
&+\mathcal X_{\alpha,\beta}(\tau)
\sum_{m,n\ge1}\frac{
 V^-_{\alpha,\beta,\tau}(mn/X)}
 {m^{\widetilde s_1}n^{\widetilde s_2}}
-\mathcal P_{\alpha,\beta}(\tau),
\end{aligned}}
\label{eq:exact-shifted-product-AFE}
\end{equation}
where
\begin{equation}
\mathcal P_{\alpha,\beta}(\tau)
:=\sum_{u_0\in\{-s_1,\,1-s_1,\,-s_2,\,1-s_2\}}
\operatorname*{Res}_{u=u_0}
\left[
 G_{\rm AFE}(u)
 \frac{\Lambda(s_1+u)\Lambda(s_2+u)}
 {\gamma(s_1)\gamma(s_2)}\frac1u
\right].
\label{eq:AFE-four-polar-residues}
\end{equation}
For every fixed $a,b,m,M\ge0$,
\begin{equation}
\boxed{
\partial_A^a\partial_B^b(T\partial_\tau)^m
\mathcal P_{\alpha,\beta}(\tau)=O_{C,a,b,m,M}(T^{-M}).}
\label{eq:AFE-polar-Gaussian-bound}
\end{equation}
The dual functional-equation multiplier is
\begin{equation}
\mathcal X_{\alpha,\beta}(\tau)
:=\frac{\gamma(1-s_1)\gamma(1-s_2)}
{\gamma(s_1)\gamma(s_2)}
=X^{-\alpha-\beta}\left(1+O_C(T^{-1})\right).
\label{eq:AFE-dual-factor}
\end{equation}

Let $p=A+B=L(\alpha+\beta)$ and put
\[
\eta_{\alpha,\beta}(\tau)
:=X^{\alpha+\beta}\mathcal X_{\alpha,\beta}(\tau)
=1+O_C(T^{-1}).
\]
Regard $U$ temporarily as an independent endpoint parameter in a fixed compact
interval containing $U_\tau$; the physical fibre is obtained by setting
$U=U_\tau$.  After the diagonal pole variable is written as
$r=\log\ell/L$, and after the dual coordinate is reflected by $v=U-r$, the
two Dirichlet-sum halves give the smoothed fibre
\begin{equation}
\boxed{
\begin{aligned}
\mathcal E^{J}_{U,T}(A,B;\tau):={}&
\int_0^\infty e^{-(A+B)r}
 V^+_{\alpha,\beta,\tau}\!\left(e^{L(2r-U)}\right)\,dr\\
&+\eta_{\alpha,\beta}(\tau)
 \int_{-\infty}^{U}e^{-(A+B)v}
 V^-_{\alpha,\beta,\tau}\!\left(e^{-L(2v-U)}\right)\,dv
-\mathcal P_{\alpha,\beta}(\tau).
\end{aligned}
}
\label{eq:J-AFE-actual-smooth-fibre}
\end{equation}
Then
\begin{equation}
\boxed{
\mathcal E^{J}_{U,T}(A,B;\tau)
=E_U(A+B)+\mathcal R^{J}_{U,T}(A,B;\tau),
\qquad
E_U(p)=\int_0^Ue^{-pr}\,dr.
}
\label{eq:J-AFE-sharp-transfer}
\end{equation}
Thus the subtracted polar term in
\eqref{eq:exact-shifted-product-AFE} is included explicitly in
$\mathcal R^{J}_{U,T}$ and satisfies the stronger bound
\eqref{eq:AFE-polar-Gaussian-bound}.  Uniformly for bounded $A,B,U$ and
$a+b+j+m\le2$,
\begin{equation}
\boxed{
\partial_A^a\partial_B^b\partial_U^j(T\partial_\tau)^m
\mathcal R^{J}_{U,T}(A,B;\tau)
=O_C\!\left(\frac{L^j}{T}+T^{-10}\right),
}
\label{eq:J-AFE-sharp-error}
\end{equation}
and hence the maximum of these derivatives is
$O_C(L^3/T+T^{-10})=o(1)$.  Along the physical relation
$\tau=2\pi e^{LU}$ one has
\[
\frac d{dU}=\partial_U+L\tau\partial_\tau;
\]
the first two total endpoint derivatives are still
$O_C(L^3/T+T^{-10})$.
The same bound holds for the one-sided diagonal derivative jump after this
fibre is inserted into any fixed finite combination of the principal profiles
used below.  Consequently the difference between the actual smoothed critical-$J$
lag form and the sharp-$E_U$ lag form is $o(1)$ in the triangular form norm
\eqref{eq:triangular-form-norm}.
\end{lemma}

\begin{proof}
Because $G_{\rm AFE}$ is even, substitute $u\mapsto-u$ in
$V(x^{-1})$ and shift the resulting line from $\Re u=-c$ to $\Re u=c$.
The only crossed pole is the simple pole at $u=0$, with residue
 $G_{\rm AFE}(0)=1$.  This gives \eqref{eq:AFE-complement}.  The same contour
shift gives, for every fixed $M,j\ge0$,
\[
(x\partial_x)^jV(x)\ll_{M,j}x^{-M}\quad(x\ge1),
\qquad
(x\partial_x)^j(1-V(x))\ll_{M,j}x^M\quad(0<x\le1).
\]

We first record the exact AFE, since this is where the completed zeta poles
matter.  For $c>\frac12+C/L$, absolute convergence gives
\[
I^+_{\alpha,\beta}(\tau)
:=\frac1{2\pi i}\int_{(c)}G_{\rm AFE}(u)
 \frac{\Lambda(s_1+u)\Lambda(s_2+u)}
 {\gamma(s_1)\gamma(s_2)}\frac{du}{u}
=\sum_{m,n\ge1}\frac{V^+_{\alpha,\beta,\tau}(mn/X)}
 {m^{s_1}n^{s_2}}.
\]
Moving the line to $\Re u=-c$ crosses $u=0$ and, because $\Lambda$ has poles
at both $0$ and $1$, precisely the four remote points
\[
-s_1,\qquad 1-s_1,\qquad -s_2,\qquad 1-s_2.
\]
Thus
\[
I^+_{\alpha,\beta}(\tau)
=\zeta(s_1)\zeta(s_2)+\mathcal P_{\alpha,\beta}(\tau)
 +I^{\rm left}_{\alpha,\beta}(\tau).
\]
Using $\Lambda(s)=\Lambda(1-s)$ on the left line and then substituting
$u\mapsto-u$ gives
\[
I^{\rm left}_{\alpha,\beta}(\tau)
=-\mathcal X_{\alpha,\beta}(\tau)I^-_{\alpha,\beta}(\tau).
\]
After interchanging $m,n$, the Dirichlet expansion of $I^-$ is the second sum
in \eqref{eq:exact-shifted-product-AFE}, with the shifts
$(-\beta,-\alpha)$.  This proves the exact identity and also explains why no
such remote residues occur in the preceding proof of the pure identity
\eqref{eq:AFE-complement}: the pure Mellin integrand contains neither
$\Lambda$ factor.

At every remote pole $u_0$ one has $|\Im u_0|=\tau+O(1)$ and
\[
|G_{\rm AFE}(u_0)|
=\exp\{-\tau^2+O(\tau/L)+O(1)\}.
\]
The gamma residues, the reciprocal denominator, and the remaining zeta factor,
together with any fixed number of $A,B,T\partial_\tau$ derivatives, contribute
at most $e^{C'\tau}(1+\tau)^{C'}$.  The Gaussian absorbs this loss and proves
\eqref{eq:AFE-polar-Gaussian-bound}.

For completeness, the normalized gamma ratio in the primal weight is
\[
R^+_{\alpha,\beta,\tau}(u)
:=X^{-u}\pi^{-u}
\frac{\Gamma((\frac12+\alpha+i\tau+u)/2)
      \Gamma((\frac12+\beta-i\tau+u)/2)}
     {\Gamma((\frac12+\alpha+i\tau)/2)
      \Gamma((\frac12+\beta-i\tau)/2)},
\]
and the dual ratio is the same expression with
$(\alpha,\beta)$ replaced by $(-\beta,-\alpha)$.  On
$|\Im u|\le T^{1/4}$, uniform Stirling expansion on fixed real strips gives
\eqref{eq:AFE-gamma-ratio} and, more generally, for
$a+b+m\le2$,
\begin{equation}
\partial_A^a\partial_B^b(T\partial_\tau)^m
\bigl(R^\pm_{\alpha,\beta,\tau}(u)-1\bigr)
\ll_C\frac{(1+|u|)^{2+a+b+m}}{T}.
\label{eq:AFE-joint-gamma-ratio}
\end{equation}
The identical expansion of the functional-equation quotient gives
\eqref{eq:AFE-dual-factor} and
\begin{equation}
\partial_A^a\partial_B^b(T\partial_\tau)^m
\bigl(\eta_{\alpha,\beta}(\tau)-1\bigr)=O_C(T^{-1})
\qquad(a+b+m\le2).
\label{eq:AFE-joint-dual-factor}
\end{equation}
Here $\partial_A=L^{-1}\partial_\alpha$ and
$\partial_B=L^{-1}\partial_\beta$.  For
$|\Im u|>T^{1/4}$, Stirling's crude polynomial growth bound combined with
$|G_{\rm AFE}(c+iy)|\ll_c e^{-y^2/2}$ contributes $O_M(T^{-M})$ for every
fixed $M$.  In particular $R^\pm-1$ vanishes at $u=0$.  Therefore the
difference $V^\pm_{\alpha,\beta,\tau}-V$, as well as every parameter
derivative in \eqref{eq:AFE-joint-gamma-ratio}, has a Mellin integrand
holomorphic at the origin.  Splitting the integral at
$|\Im u|=T^{1/4}$ and shifting its central part to
$\Re u=\pm c_0$ gives, for $a+b+m+q\le2$,
\begin{equation}
\left|
\partial_A^a\partial_B^b(T\partial_\tau)^m(x\partial_x)^q
\bigl(V^\pm_{\alpha,\beta,\tau}(x)-V(x)\bigr)\right|
\ll_C T^{-1}\min(x^{c_0},x^{-c_0})+O_M(T^{-M}).
\label{eq:AFE-weight-difference}
\end{equation}
The horizontal connectors introduced by this finite-height truncation are
$O_M(T^{-M})$ by the same Gaussian bound.

At diagonal level write $r=\log\ell/L$.  The primal shift factor is $e^{-pr}$,
and the pure primal weight is
\[
V\!\left(e^{L(2r-U)}\right),
\qquad r\ge0.
\]
In the dual half the shift is reversed, so before reflection its scalar factor
is $\mathcal X_{\alpha,\beta}(\tau)e^{pr}$.  With $v=U-r$ this becomes
\[
\mathcal X_{\alpha,\beta}(\tau)e^{p(U-v)}
=\eta_{\alpha,\beta}(\tau)e^{-pv},
\]
and the pure dual weight becomes
\[
V\!\left(e^{-L(2v-U)}\right),
\qquad v\le U.
\]
Thus, after replacing the gamma ratios by their pure Mellin parts and
$\eta_{\alpha,\beta}$ by $1$, the two scalar factors are the same $e^{-p\cdot}$
on the reflected fibre.  On the overlap $0\le v\le U$,
\eqref{eq:AFE-complement} says that the two weights add \emph{exactly} to one.
Thus the pure smoothed halves give exactly $E_U(p)$ on the full fibre.  The only
pure-weight pieces outside that interval are $r>U$ and $v<0$; there the Mellin
argument is at least $e^{LU}=X\asymp T$, so the preceding rapid-decay estimate
makes their contribution, together with all the stated joint derivatives,
$O(T^{-10})$ after choosing the auxiliary decay order $M$ large enough.

It remains to replace the exact gamma weights by the pure weight.  With
$x=e^{L(2r-U)}$, \eqref{eq:AFE-weight-difference} is exponentially localized to
a strip of normalized width $O(L^{-1})$ about $r=U/2$.  Each $U$ derivative
costs at most one factor $L$, while $A,B$ derivatives cost only bounded powers
of $r$ on the retained fibre.  Integrating the localized bound and using
\eqref{eq:AFE-joint-dual-factor}, then adding the Gaussian-small contribution
\eqref{eq:AFE-polar-Gaussian-bound}, gives the sharper
$O_C(L^j/T+T^{-10})$ bound and hence
\eqref{eq:J-AFE-sharp-error}.  Replacing $\partial_U$ by
$\partial_U+L\tau\partial_\tau$ for the physical endpoint costs at most $L^2$
through second order, which is covered by the displayed conservative
$O(L^3/T)$ bound.  The same calculation gives the corresponding diagonal jump
and regular second-lag derivative bounds.  Applying the two-variable primitive identity
used in Section~\ref{sec:model-floor} then converts these bounds into an
$o(1)$ triangular-form perturbation.  This proves the lemma.
\end{proof}

\begin{proposition}[The critical-$J$ self diagonal transfers to the same full fibre]
\label{prop:J-full-fibre}
At the pointwise principal diagonal, the actual smoothed primal and dual
approximate-functional-equation halves of the critical-$J$ self source,
together with the subtracted four-residue polar correction in
\eqref{eq:exact-shifted-product-AFE}, equal a sharp full-fibre profile plus the
triangularly negligible remainder of
Lemma~\ref{lem:J-AFE-sharp-transfer}.  In particular, if
\[
U:=\frac{\log X}{L},
\qquad X=\frac{t}{2\pi},
\]
and $A:=L\alpha$, $B:=L\beta$, then the sharp principal representatives of the
two $h=k=1$ diagonal halves are
\[
E_{U/2}(A+B)
\]
and
\[
e^{-U(A+B)}E_{U/2}(-(A+B)),
\]
and their sum is exactly
\begin{equation}
\boxed{
E_{U/2}(A+B)
+e^{-U(A+B)}E_{U/2}(-(A+B))
=E_U(A+B).
}
\label{eq:J-half-glue}
\end{equation}
The difference between the actual smoothed $J$ fibre and this sharp profile is
$o(1)$ in the triangular form norm and is assigned to the regular remainder,
not to $P_T$.

For the completed three-$Z$ diagonal, Lemma~\ref{lem:HLZ-local-height-rescaling}
permits the fixed HLZ Mellin scale to be replaced by the same local height
$X=t/(2\pi)$ without changing the $u=0$ main residue.  After
$m=k\ell,n=h\ell$ its principal radial factor is therefore
\[
\left(\frac{X^2}{hk\ell^2}\right)^u,
\]
so the three-$Z$ principal term uses the same full logarithmic endpoint.  The
factors $h,k$ do not create a separate principal endpoint after the common
analysis map is fixed: by Lemma~\ref{lem:outer-translation-intertwiner} their
logarithms are absorbed by exact packet translations, leaving the universal
packet lag.  Thus neither a smooth-to-sharp identification nor a global-to-local
Mellin-scale substitution is being made without an explicit error estimate.
\end{proposition}

\begin{proof}
For the shifted product
\[
\zeta(1/2+it+\alpha)\zeta(1/2-it+\beta),
\]
use the exact identity \eqref{eq:exact-shifted-product-AFE}, whose two
Dirichlet sums have the weights \eqref{eq:exact-AFE-weights}.  On the diagonal
the radial conductor is
$mn/X$, and with $h=k=1$ one has $mn=\ell^2$.  Lemma~\ref{lem:J-AFE-sharp-transfer}
shows directly, using the complement identity \eqref{eq:AFE-complement}, that
after reflecting the dual coordinate the two smooth diagonal fibres and the
subtracted polar contribution give
\[
E_U(A+B)+\mathcal R^J_{U,T}(A,B;t),
\]
with the derivative and triangular-form bounds
\eqref{eq:J-AFE-sharp-error}.  The sharp part may be split at $U/2$ into the two
displayed terms, and the elementary change of variables $v=U-r$ gives
\eqref{eq:J-half-glue}.  This proves the critical-$J$ assertion without
identifying a smooth AFE cutoff with a characteristic function.

For the three-$Z$ term, begin with the actual HLZ diagonal factor
$T^{2u}(hk\ell^2)^{-u}$.  Apply
Lemma~\ref{lem:HLZ-local-height-rescaling} at the physical height $t$; the
multiplier $(X/T)^{2u}$ belongs to the same admissible Mellin class, is equal to
one at $u=0$, and hence leaves the main residue unchanged.  The resulting
principal factor is exactly $(X^2/(hk\ell^2))^u$, which yields the same full
logarithmic interval $0\le\log\ell/L\le U$.  The remaining $h,k$ dependence is
removed by the exact unitary translations of
Lemma~\ref{lem:outer-translation-intertwiner}.  This completes the comparison.
\end{proof}

For the remainder of the proof fix the stable logarithmic core
\begin{equation}
\boxed{
W_T:=L-A_0\log L,
\qquad
X_T:=e^{W_T}=\frac{T}{2\pi L^{A_0}},
}
\label{eq:stable-log-core}
\end{equation}
where the fixed lower bounds on $A_0$ are imposed in
Section~\ref{sec:retained}. This definition belongs to the invariant-source
geometry and is independent of the auxiliary divisor comparison below.

\begin{proposition}[Auxiliary unsummed common-divisor feature]
\label{prop:common-divisor-feature}
For comparison with the general twisted formula, let $d,e\le Y$ be auxiliary
outer Dirichlet-polynomial indices, with coefficients
$c_d,c_e$, and let $m,n$ be the two zeta indices. For the invariant contour
source Proposition~\ref{prop:trivial-outer-twist} specializes this family to
$d=e=1$; the general finite-range identity below is auxiliary. On the stationary diagonal
\begin{equation}
\boxed{dm=en=:N,}
\label{eq:common-product-N}
\end{equation}
the base principal shifted product has the exact positive factorization
\begin{equation}
\boxed{
\sum_NN^{-1-2a}
\left|\sum_{\substack{d\mid N\\d\le Y}}c_dd^a\right|^2,
\qquad a=\frac pL.
}
\label{eq:common-divisor-gram}
\end{equation}
The translated principal master with $\delta=2/L$, inner shifts
$a+\delta,a-\delta$, and reciprocal outer deformation
$c_d\mapsto c_dd^{-\delta}$, $c_e\mapsto c_ee^{\delta}$ has exactly the same
factor \eqref{eq:common-divisor-gram}. The critical-$J$ self slot acts through
the same outer-divisor feature. The factorization is Hilbert-valued as well: if
$b_d\in\mathcal H$ and $U_d$ are unitaries on $\mathcal H$, then the corresponding
outer feature is
\begin{equation}
\boxed{
\sum_NN^{-1-2a}
\left\|\sum_{\substack{d\mid N\\d\le Y}}d^aU_db_d\right\|_{\mathcal H}^2
\ge0.
}
\label{eq:common-divisor-Hilbert-gram}
\end{equation}
Thus the exact packet translations, including the full smooth packet profile,
may be retained inside the common positive analysis map. Hence, before scalar
contraction, the completed signed principal dependence is radial and the outer
arithmetic factor is one common positive congruence.

If $d=gh$, $e=gk$, $(h,k)=1$, and
$m=k\ell$, $n=h\ell$, then
\begin{equation}
\boxed{\ell=\frac{Ng}{de}.}
\label{eq:ell-N-displacement}
\end{equation}
Consequently, if this auxiliary finite family is parametrized by
$Y=L^\kappa$ for a fixed $\kappa>0$, then
\begin{equation}
0\le\frac{\log N-\log\ell}{L}
=\frac{\log(de/g)}{L}
\ll_\kappa\frac{\log L}{L}.
\label{eq:outer-endpoint-displacement}
\end{equation}
The same $O(\log L/L)$ scale controls the $p$-dependence of the positive outer
divisor map. Moreover $\mathcal W_a$ itself has an explicit finite M\"obius
inverse on the outer range: restricting its output to $N\le Y$,
\begin{equation}
\boxed{
 c_dd^a=\sum_{e\mid d}\mu(d/e)(\mathcal W_ac)(e),
 \qquad d\le Y.
}
\label{eq:Wa-mobius}
\end{equation}
Consequently the weighted map
$c\mapsto(N^{-1/2-a}(\mathcal W_ac)(N))_{N\le Y}$ is injective and has inverse
norm $Y^{1+o(1)}$. The same inverse bound holds for the Hilbert-valued map in
\eqref{eq:common-divisor-Hilbert-gram}: if
$S_N=\sum_{d\mid N}d^aU_db_d$, then
\[
d^aU_db_d=\sum_{e\mid d}\mu(d/e)S_e,
\]
so $b_d$ is recovered by applying the unitary $U_d^*$ and the scalar weight
$d^{-a}$. The deformation of this positive divisor map has entrywise
normalized size $O(\log L/L)$. To avoid relying on one-sided factorization in
any later comparison, we measure it with the uniform two-sided arithmetic
transfer bound of Lemma~\ref{lem:arithmetic-loss}; hence
\begin{equation}
\boxed{O\!\left(L^{-1+10\kappa+o(1)}\log L\right)=o(1),
\qquad \kappa<\frac1{10}.}
\label{eq:outer-divisor-deformation}
\end{equation}
For the invariant contour source, $d=e=1$ and hence the divisor-map
$p$-deformation, the endpoint displacement
\eqref{eq:outer-endpoint-displacement}, and the outer annulus all vanish
identically. For the auxiliary general finite family, the endpoint displacement
is pair dependent and extends the stationary endpoint in the common
$N$-coordinate; Proposition~\ref{prop:outer-excess-edge} records its separate
sideband decomposition. None of these auxiliary losses is needed for the exact
source specialization.
\end{proposition}

\begin{proof}
A base diagonal monomial is
\[
c_d\overline{c_e}\,
d^{-s}e^{-(1-s)}m^{-(s+a)}n^{-(1-s+a)}.
\]
Using $dm=en=N$, equivalently $m=N/d$ and $n=N/e$, this is exactly
\[
N^{-1-2a}c_dd^a\,\overline{c_ee^a}.
\]
Summing first over divisors of $N$ proves \eqref{eq:common-divisor-gram}. More
generally, replacing the scalar factors by
$x_d=d^aU_db_d\in\mathcal H$ gives
\[
\sum_{d,e\mid N}\langle x_d,x_e\rangle
=\left\|\sum_{d\mid N}x_d\right\|^2,
\]
which proves \eqref{eq:common-divisor-Hilbert-gram}. For the translated master
the extra scalar factor is
\[
d^{-\delta}e^{\delta}m^{-a-\delta}n^{-a+\delta}
=N^{-2a}d^ae^a,
\]
so the same divisor Gram is obtained identically. The critical-$J$ source
multiplies the same shifted product, and Proposition~\ref{prop:J-full-fibre}
shows that its primal and dual halves glue before scalar contraction; therefore
its outer analysis map is the same one.

For the inverse statement, write
$S_N=(\mathcal W_ac)(N)$. Then
$S_N=\sum_{d\mid N}c_dd^a$ for $N\le Y$, and ordinary M\"obius inversion gives
\eqref{eq:Wa-mobius}. Cauchy--Schwarz and the divisor bound
$\tau(n)=n^{o(1)}$ give
\[
\sum_{d\le Y}|c_d|^2
\ll Y^{1+o(1)}\sum_{N\le Y}|S_N|^2
\ll Y^{2+o(1)}\sum_{N\le Y}N^{-1-2a}|S_N|^2,
\]
which is the asserted $Y^{1+o(1)}$ inverse norm.

Finally $N=ghk\ell$ and $de=g^2hk$, proving
\eqref{eq:ell-N-displacement}. Since $d,e\le Y$, one has
$de/g\le Y^2$, giving \eqref{eq:outer-endpoint-displacement}. In particular,
the physical cutoff
\[
\ell\le Y_{\rm fib}e^{-|v|}
\]
is equivalent in the common-divisor coordinate to
\begin{equation}
N\le \frac{de}{g}Y_{\rm fib}e^{-|v|}.
\label{eq:N-actual-endpoint}
\end{equation}
Thus the pair-independent range
\begin{equation}
N\le Y_{\rm fib}e^{-|v|}
\label{eq:N-common-core}
\end{equation}
is a genuine common sub-prefix for every arithmetic pair; the difference
between \eqref{eq:N-actual-endpoint} and \eqref{eq:N-common-core} is an outer
annulus, not a shortening of the common core. After applying
Lemma~\ref{lem:outer-translation-intertwiner}, the raw lag $v$ is straightened
to the universal packet lag $w$. Imposing the common core then simply truncates
the outer sum in \eqref{eq:common-divisor-gram}:
\begin{equation}
\boxed{
\sum_{N\le Y_{\rm fib}e^{-|w|}}
N^{-1-2a}
\left|\sum_{\substack{d\mid N\\d\le Y}}c_dd^a\right|^2\ge0.
}
\label{eq:common-core-positive-gram}
\end{equation}
Thus the passage to the pair-independent core preserves the exact positive
arithmetic congruence; only the discarded annulus is treated separately.
The same truncation also preserves the coordinates needed for the finite
M\"obius inverse. Indeed, on this stable core,
\[
|w|\le W_T,
\]
while $\log Y_{\rm fib}=L+O(1)$. Hence uniformly on the retained core
\begin{equation}
Y_{\rm fib}e^{-|w|}
\ge L^{A_0+o(1)}.
\label{eq:common-core-inverse-range}
\end{equation}
For this auxiliary comparison one may, after $A_0$ is fixed, choose any
$0<\kappa<\min\{1/10,A_0-1\}$ and set $Y=L^\kappa$. Then the right side exceeds
$Y$ for all sufficiently large $T$. Consequently every coordinate
$1\le N\le Y$ used in \eqref{eq:Wa-mobius} remains present after the common-core
restriction, and the inverse bounds for $\mathcal W_a$ and
$\widetilde{\mathcal W}_a$ remain valid. Moreover
$d^a=\exp(p\log d/L)$, so each fixed $p$-derivative of the positive divisor
map costs at most a fixed power of $\log Y/L$. Applying the uniform two-sided
transfer estimate \eqref{eq:arithmetic-two-sided-transfer} gives
\eqref{eq:outer-divisor-deformation}. This optional auxiliary choice is not part
of the invariant-source parameter hierarchy.
\end{proof}

\begin{proposition}[Auxiliary outer common-prefix excess]
\label{prop:outer-excess-edge}
For the invariant contour source this contribution is identically zero by
Proposition~\ref{prop:trivial-outer-twist}. For a general auxiliary retained
outer pair $(d,e)$, first restrict to the stable logarithmic
band $0\le u\le W_T$. On this band decompose the actual $N$-range
\eqref{eq:N-actual-endpoint} into the common core \eqref{eq:N-common-core} and
the excess annulus
\begin{equation}
Y_{\rm fib}e^{-|v|}<N
\le \frac{de}{g}Y_{\rm fib}e^{-|v|}.
\label{eq:outer-excess-annulus}
\end{equation}
Let $A_T^{\rm out}$ denote this stable-band annulus and its adjoint; the part on
$W_T<u\le L$ belongs only to \eqref{eq:terminal-log-tail}. After returning to
the physical stationary carrier $\ell=Ng/(de)$, expand the smooth unit-shell
stationary symbol in Fourier sidebands and write
\begin{equation}
\boxed{A_T^{\rm out}=E_T^{\rm out}+R_T^{\rm out,hi},}
\label{eq:outer-excess-sideband-split}
\end{equation}
where $E_T^{\rm out}$ contains the sidebands
$|r|\le R_j^{\rm sb}$ and $R_T^{\rm out,hi}$ the complementary sidebands.
Then
\begin{equation}
\boxed{
\rank E_T^{\rm out}
\ll T Y^{3+o(1)}=o(N_T).
}
\label{eq:outer-excess-edge-rank}
\end{equation}
The term $R_T^{\rm out,hi}$ belongs to the direct high-sideband regular
class~(vii) of Proposition~\ref{prop:regular-ledger}. No estimate of the annulus
width is needed for the low-sideband rank bound.
\end{proposition}

\begin{proof}
Equation~\eqref{eq:ell-N-displacement} maps the annulus back to positive integer
carriers $\ell$ satisfying the original stationary saddle relation
$t_*=2\pi\ell y$. If $y\in[J_j,J_j+1]$ and $t_*\in[T,2T]$, then
\[
\frac{T}{2\pi(J_j+1)}\le \ell\le \frac{2T}{2\pi J_j},
\]
which is precisely the stationary carrier interval used in
Section~\ref{sec:edge-rank}, up to the fixed shell buffer already present there.
Lemma~\ref{lem:stationary-symbol} gives a smooth symbol uniformly on the
annulus. Its low Fourier sidebands replace $\ell$ by effective shifts within
$R_j^{\rm sb}$ of this stationary interval. Proposition~\ref{prop:edge-rank}
therefore places all low-sideband terms in the same shell-edge spaces and gives
total geometric rank $O(T)$ per fixed arithmetic multiplicity. Multiplying by
the conservative $Y^{3+o(1)}$ outer label count proves
\eqref{eq:outer-excess-edge-rank}. The high sidebands are not included in this
rank assertion; their coefficient square sum is no larger than that of the full
direct family and their Fourier decay is estimated in class~(vii).
\end{proof}

For a packet pair $r,s$, write $w_r,w_s\in\mathbb R$ for their logarithmic
packet displacements in the stationary dilation of
\eqref{eq:stationary-dilation}. Let $I_r$ and $I_s$ be the corresponding two
physical logarithmic prefix intervals and put
\[
J_{rs}:=I_r\cap I_s,
\qquad
\rho_{rs}^{(T)}:=|J_{rs}|.
\]
In the invariant-source specialization there is no outer translation and
$N=\ell$. Propositions~\ref{prop:J-full-fibre} and \ref{prop:X-tagging} show
that the principal full-fibre scale is the local archimedean height
$Y_{\rm fib}=X_{\rm arch}=\tau/(2\pi)$. Hence
\[
\nu_T:=\frac{\log Y_{\rm fib}}{L}=\lambda(\tau)
=1+O(L^{-1}),
\qquad \nu_T\ge1
\]
on the dyadic interval.
The actual normalized common-prefix endpoint is
\begin{equation}
\boxed{
\rho_{rs}^{(T)}=\nu_T-\frac{|w_r-w_s|}{L}.
}
\label{eq:finiteT-common-endpoint}
\end{equation}
The frozen global model replaces the scalar full-fibre length $\nu_T$ by $1$.
This $O(L^{-1})$ endpoint deformation is one component of the parameter defect
\eqref{eq:finiteT-parameter-defect} and is estimated in
Lemma~\ref{lem:finiteT-parameter-regular}, not by a variable-parameter
positivity transfer.  In the finite-$T$ local formulas below we write
$\rho_{rs}$ for $\rho_{rs}^{(T)}$. Since $d=e=g=1$ in the invariant source,
there is no additional arithmetic endpoint displacement.
The intersection $I_r\cap I_s$ is formed at the packet level before the local
Perron residue is evaluated. Thus $\rho_{rs}$ is determined by the global
packet geometry rather than introduced as an independent positivity parameter.

Before the $u$-residue and the three shift residues are taken, the packet pair
contributes one joint analytic weight on the pair-independent common $N$-core.
Here $t$ denotes the normalized logarithmic $N$-coordinate after translating
that core interval to its local origin; it is not the physical height variable.
Thus
\begin{equation}
W_{rs,L}(u,\mathbf z)
=\int_{J_{rs}}a_{rs}(t,\mathbf z)e^{Lut}\,dt,
\label{eq:joint-weight}
\end{equation}
where $a_{rs}$ contains the smooth packet polarization together with the slow
archimedean and fixed shift factors.  The dependence on $r,s$ is entirely the
packet-pair dependence made explicit in
Lemma~\ref{lem:joint-weight-common-factor} below.
For every fixed multi-index $\nu$,
\begin{equation}
\left|\partial_{\mathbf z}^{\nu}[u^n]W_{rs,L}(u,\mathbf z)\right|
\ll_\nu \frac{(CL)^n(1+n)^C}{n!}.
\label{eq:joint-factorial}
\end{equation}

\begin{lemma}[Shift-independent common packet analysis for the finite model]
\label{lem:joint-weight-common-factor}
For the finite model source $H_{\rm mod}$, there is a Hilbert output space
$\mathcal H_{\rm fib}$ and one linear packet map
\[
\mathscr C_T:\mathbb C^{\mathcal I_T}\longrightarrow\mathcal H_{\rm fib},
\]
independent of the auxiliary Cauchy shift vector $\mathbf z$.  For every local
Mellin variable $u$ and every $\mathbf z\in\mathcal T_L$ there is a
packet-label-independent multiplication operator $M_{u,\mathbf z}$ such that
\begin{equation}
\boxed{
W_{rs,L}(u,\mathbf z)
=\left\langle
M_{u,\mathbf z}\mathscr C_Te_s,
\mathscr C_Te_r
\right\rangle_{\mathcal H_{\rm fib}}.
}
\label{eq:joint-weight-matrix-coefficient}
\end{equation}
More explicitly, the principal one-sided packet/stationary factor may be chosen
as a function $b_{r,T}(t)$ independent of $\mathbf z$, while all auxiliary
shift dependence is retained in a common scalar factor $c_T(t,\mathbf z)$:
\[
\mathbf 1_{J_{rs}}(t)a_{rs}(t,\mathbf z)
=
\mathbf 1_{I_s}(t)b_{s,T}(t)\,
\overline{\mathbf 1_{I_r}(t)b_{r,T}(t)}\,
c_T(t,\mathbf z).
\]
Thus one may take
\[
(\mathscr C_Te_r)(t):=\mathbf 1_{I_r}(t)b_{r,T}(t),
\qquad
(M_{u,\mathbf z}F)(t):=c_T(t,\mathbf z)e^{Lut}F(t).
\]
Consequently every finite Cauchy derivative acts on the output operator and
never produces jets of the packet analysis map.
\end{lemma}

\begin{proof}
The packet polarization \eqref{eq:packet-polarization} contains no HLZ Cauchy
shift.  At principal stationary level, the exact dilation
\eqref{eq:dilation-covariance} acts only on the one-sided packet factor and its
support interval.  The completed three-$Z$ factors, the archimedean factors,
the Cauchy denominators, and the Mellin factor depend on the auxiliary shifts
but not on the packet label. Lemma~\ref{lem:model-HLZ-Fubini-lift} shows that
for the finite model source the scalar HLZ transfer preserves this separation:
linearity and Fubini leave the one-sided packet factors untouched. Hence the
shift-dependent common factors may be placed entirely in $c_T(t,\mathbf z)$,
leaving $b_{r,T}(t)$ independent of $\mathbf z$.  Since
$I_r\cap I_s=J_{rs}$, substitution gives
\eqref{eq:joint-weight-matrix-coefficient}.  Differentiating the finite Cauchy
representation therefore differentiates $M_{u,\mathbf z}$ alone; terms such as
$(\partial_{\mathbf z}\mathscr C_T)^*M\mathscr C_T$ do not arise.
\end{proof}

If
\[
\mathcal F(u,\mathbf z)=\sum_{j=-J}^\infty F_j(\mathbf z)u^j,
\]
then
\[
\Res_{u=0}(W\mathcal F)
=\sum_{n=0}^{J-1}[u^n]W\,F_{-n-1}.
\]
This is a finite sum, hence the $u$-residue and the finite three-$Z$ shift
residues commute exactly:
\begin{equation}
\boxed{
\Res_{\mathbf z}\Res_{u=0}(W\mathcal F)
=\Res_{u=0}\Res_{\mathbf z}(W\mathcal F).
}
\label{eq:residue-commute}
\end{equation}
On the separated Cauchy torus, the explicit factor
$H_{\mathbf z}$ in \eqref{eq:HLZ-explicit-auxiliary} cancels the moving
divisors $2u=z_2-z_1$ and $2u=z_2-z_3$ pointwise; multiplication by the
holomorphic packet factor \eqref{eq:joint-weight} does not recreate them.

For the basic local pole feature,
\[
\int_{J_{rs}}e^{-pt}\,dt
\]
becomes, after translation of $J_{rs}$ to $[0,\rho_{rs}]$,
\[
E_{\rho_{rs}}(p):=\int_0^{\rho_{rs}}e^{-pt}\,dt.
\]
The following proposition records the operator statement needed to pass from
this scalar profile to the actual packet matrix.

\begin{proposition}[Common-core factorization and frozen principal split]
\label{prop:common-core-congruence}
Restrict the finite model packet source to the stable logarithmic core
$0\le u\le W_T$.  Choose a nonnegative bounded-overlap height partition
\begin{equation}
\boxed{
\sum_b\omega_b(\tau)=1,
\qquad
\operatorname{supp}\omega_b\subset\{\,|\tau-\tau_b|\ll H\,\},
\qquad H=T/L^D,
}
\label{eq:height-principal-partition}
\end{equation}
with
\begin{equation}
\boxed{
|\partial_\tau^j\omega_b(\tau)|\ll_jH^{-j}
\qquad(0\le j\le3).
}
\label{eq:height-partition-derivatives}
\end{equation}
The partition is inserted \emph{linearly} in the physical-height contour
integral before stationary phase.  On the flat packet plateau the common
one-sided carrier analysis is
\begin{equation}
\boxed{
(\mathcal A_T^{\rm core}v)(u)
:=L_{\rm pkt}^{-1/2}\sum_{\nu\in\mathcal I_T}
 e^{i\tau_\nu u_0}v_\nu e^{i\tau_\nu u},
\qquad 0\le u\le W_T,
}
\label{eq:common-core-analysis}
\end{equation}
where $u_0$ is the fixed translation of the physical Fourier plateau.  Put
$R_T^{\rm core}=W_T/L$ and define
\begin{equation}
\boxed{
(\widehat{\mathcal A}_T^{\rm core}v)(r)
:=L^{1/2}(\mathcal A_T^{\rm core}v)(Lr),
\qquad 0\le r\le R_T^{\rm core}.
}
\label{eq:normalized-common-core-analysis}
\end{equation}
After the finite Cauchy residues, the local-height HLZ rescaling of
Lemma~\ref{lem:HLZ-local-height-rescaling}, and the sharp-fibre transfer of
Lemma~\ref{lem:J-AFE-sharp-transfer}, all three completed model faces have this
same one-sided packet/core analysis.

For a packet pair let
\[
\rho_{rs}^{(0)}:=1-\frac{|w_r-w_s|}{L},
\qquad
\delta(\tau):=\lambda(\tau)-1.
\]
Then \eqref{eq:finiteT-common-endpoint} gives
$\rho_{rs}^{(T)}=\rho_{rs}^{(0)}+\delta(\tau)$.  Define the scalar finite-$T$
parameter defect, before stationary shell contraction, by
\begin{equation}
\boxed{
\Delta_{b,\tau}(\rho;p)
:=\mathscr G_{\lambda(\tau),\mu_b,\rho+\delta(\tau)}(p)
 -\mathscr G_{1,1,\rho}(p).
}
\label{eq:finiteT-parameter-defect}
\end{equation}
Uniformly for $0\le\rho\le1$, bounded $p$, and $0\le j,k\le2$,
\begin{equation}
\boxed{
|\partial_\rho^j\partial_p^k\Delta_{b,\tau}(\rho;p)|\ll L^{-1}.
}
\label{eq:finiteT-parameter-defect-derivatives}
\end{equation}
Consequently the sharp finite-model main on the stable core has the exact
algebraic split
\begin{equation}
\boxed{
A_{T,\rm mod}^{\rm sharp}=P_T^{\rm fr}+R_T^{\rm par},
\qquad
P_T^{\rm fr}
=(\widehat{\mathcal A}_T^{\rm core})^*K_G
 \widehat{\mathcal A}_T^{\rm core},
}
\label{eq:common-core-principal-congruence}
\end{equation}
where $K_G(r,s)=G(1-|r-s|)$ and $R_T^{\rm par}$ is obtained by retaining
\eqref{eq:finiteT-parameter-defect} on its original physical-height and
stationary-carrier channel.  In particular, no claim is made that
$R_T^{\rm par}$ is a scalar variable-coefficient lag kernel in the final
$(r,s)$ variables.  Its relative Hilbert--Schmidt estimate is proved in
Lemma~\ref{lem:finiteT-parameter-regular}.
\end{proposition}

\begin{proof}
By Lemma~\ref{lem:joint-weight-common-factor}, before the finite shift residues
the packet pair occurs as two copies of one one-sided analysis, while the
completed arithmetic, archimedean, Cauchy, and Mellin factors act as common
scalars on the output.  The height partition
\eqref{eq:height-principal-partition} is inserted in the original contour
integral, which is linear in the physical-height integrand.  Hence overlapping
blocks contribute the sum $\sum_b\omega_b(\tau)(\cdots)$; no term containing a
product of two different block weights is ever created.

On the stable core the Fourier variables lie strictly inside the flat plateau
of $\phi_L$, so the packet-coordinate part of the leading one-sided map is
exactly \eqref{eq:common-core-analysis}.  The principal stationary map carries
the packet factor and its support by the exact dilation
\eqref{eq:dilation-covariance}.  The exterior reserve and transition caps are
split into their low-rank stationary and relative-regular parts by
Proposition~\ref{prop:exterior-packet-support}; differentiated height cutoffs
and higher stationary/Stirling coefficients are separated into the regular
ledger. None of these pieces is part of the frozen principal term.

Proposition~\ref{prop:J-full-fibre} identifies the sharp critical-$J$ fibre
with the same full interval as the three-$Z$ faces, while
Lemma~\ref{lem:HLZ-local-height-rescaling} makes the three-$Z$ diagonal use the
same local physical scale $X=\tau/(2\pi)$.  Thus at fixed physical height and
block the scalar sharp completed profile is precisely
\eqref{eq:two-parameter-profile} with endpoint
$\rho_{rs}^{(T)}$.  Subtract and add the profile with
$\lambda=\mu=1$ and endpoint $\rho_{rs}^{(0)}$; this is the exact identity
\eqref{eq:finiteT-parameter-defect} and gives the first equality in
\eqref{eq:common-core-principal-congruence} after the common packet factors are
restored.

It remains only to verify the uniform size of the defect.  By
Lemma~\ref{lem:principal-parameter-symbol},
\[
|\lambda(\tau)-1|+|\mu_b-1|+|\delta(\tau)|\ll L^{-1}.
\]
For bounded $p$, the function
$\mathscr G_{\lambda,\mu,\rho}(p)$ of
\eqref{eq:two-parameter-profile} is analytic in $(\lambda,\mu,\rho,p)$ on a
fixed neighbourhood of
$\{1\}\times\{1\}\times[0,1]\times\{|p|\le P\}$.  The mean-value theorem,
after any fixed $\rho$- and $p$-derivatives through order two, therefore gives
\eqref{eq:finiteT-parameter-defect-derivatives}.  The frozen summand is
packet-label independent after the common one-sided analysis and hence is
exactly the pullback of $K_G$ through
\eqref{eq:normalized-common-core-analysis}.  We deliberately leave the defect
on the physical-height channel rather than converting its height dependence
into an unproved two-variable scalar coefficient field.  This proves the
proposition.
\end{proof}

\begin{lemma}[Independent face audit for the finite-model profile]
\label{lem:finite-model-face-audit}
Let $X=\tau/(2\pi)=e^{\lambda L}$, put $a=p/L$ and
$\delta=2/L$, and work at the principal pole level before the normalized
curvature coefficient is extracted.  After the two Hermitian right-edge
orientations are reassembled and the common positive $L^2$ local normalization
is removed, the three finite-model sources
\[
2F_b,\qquad -2P(s),\qquad P(s-\delta)
\]
contribute respectively
\begin{equation}
\boxed{
\mu_b e^{\lambda p}E_\rho(2p),\qquad
-2e^{\lambda p}E_\rho(p)^2,\qquad
e^{\lambda p}e^{2\lambda}E_\rho(p+2)E_\rho(p-2),
}
\label{eq:independent-three-face-audit}
\end{equation}
where $\mu_b=2F_b/L$.  In particular the two products in
\eqref{eq:independent-three-face-audit} arise from the two one-sided pole/packet
legs of one common HLZ diagonal and not from two independent arithmetic radial
variables.  The explicit factors $e^{\lambda p}$ and $e^{2\lambda}$ remain
archimedean height factors when the packet endpoint is shortened from
$\lambda$ to $\rho$.
\end{lemma}

\begin{proof}
We recompute the three coefficients from the original HLZ local factors rather
than from Proposition~\ref{prop:completed-local-profile}.  First the Hardy gauge
fixes the physical shifts.  Since $q(s)^2=\chi(1-s)$ and
$\zeta(s-a)=\chi(s-a)\zeta(1-s+a)$,
\begin{equation}
\mathscr Z(s+a)\mathscr Z(s-a)
=\frac{q(s+a)}{q(s-a)}\,
 \zeta(s+a)\zeta(1-s+a).
\label{eq:Hardy-HLZ-shift-identification}
\end{equation}
Thus the HLZ shifts are $\alpha=\beta=a$.  Uniform Stirling expansion on the
retained height block gives for the principal part
\[
\frac{q(s+a)}{q(s-a)}=X^a(1+\text{regular error}),
\]
so the common principal Hardy factor is
$X^a=e^{\lambda p}$.  The discarded correction belongs to the already-recorded
finite-height/archimedean regular class and plays no role in the principal
profile.

For the self source use the $h=k=1$ pole part of the HLZ $J$ bracket.  With
$\alpha=\beta=a$ it is
\[
\zeta(1+2a)+X^{-2a}\zeta(1-2a).
\]
Replacing the two zeta functions by their principal simple poles and including
the common Hardy factor gives
\begin{align*}
X^a\left(\frac1{2a}-\frac{X^{-2a}}{2a}\right)
&=L e^{\lambda p}E_\lambda(2p).
\end{align*}
The full Hermitian self source is $2F_b$.  After division by the common $L^2$
local normalization this is
\[
\frac{2F_b}{L}e^{\lambda p}E_\lambda(2p)
=\mu_b e^{\lambda p}E_\lambda(2p).
\]

Next consider one unit logarithmic-derivative source $P=-\zeta'/\zeta$.
For $h=k=1$ the pole part of the original HLZ factor is
\[
Z^{\rm pole}_{A,B,C}=\frac{B}{(A+B)(B+C)}.
\]
The three completed faces at $A=B=a$ therefore have pole sum
\begin{equation}
F_a(\gamma)
=\frac1{2(a+\gamma)}
 -\frac{X^{-2a}}{2(a-\gamma)}
 +X^{-a-\gamma}\frac{\gamma}{a^2-\gamma^2}.
\label{eq:HLZ-three-face-pole-sum-audit}
\end{equation}
The convention $P=-\zeta'/\zeta$ corresponds to $-\partial_\gamma$, and direct
differentiation gives
\begin{equation}
\boxed{
-\partial_\gamma F_a(0)
=\frac{(1-X^{-a})^2}{2a^2}
=\frac{L^2}{2}E_\lambda(p)^2.
}
\label{eq:HLZ-one-orientation-P-audit}
\end{equation}
This is one right-edge orientation.  The adjoint orientation gives the same
principal real scalar, so after the two orientations are reassembled and the
common $L^2$ normalization is removed, one unit $P$ source contributes
\[
e^{\lambda p}E_\lambda(p)^2.
\]
The coefficient $-2$ in $H_{\rm mod}$ therefore gives exactly
$-2e^{\lambda p}E_\lambda(p)^2$; the factor $1/2$ in
\eqref{eq:HLZ-one-orientation-P-audit} is neither missing nor duplicated.

Finally translate the source by $\delta=2/L$.  In the HLZ diagonal the exact
source-level tag is, with $h=k=1$,
\[
X^\delta(n_2n_3)^{-\delta}.
\]
On the unshifted logarithmic face the two one-sided factors carry
$n_2^{-a}$ and $n_3^{a}$ before the second orientation is placed in the common
positive fibre coordinate.  Hence the translation changes them to
\[
n_2^{-(a+\delta)},\qquad n_3^{a-\delta},
\]
which changes the two positive fibre exponents from $(p,p)$ to
$(p+2,p-2)$.  The explicit height tag is
$X^\delta=e^{2\lambda}$, while the simple-pole residue and the
orientation factor in \eqref{eq:HLZ-one-orientation-P-audit} are unchanged.
Thus the translated unit source contributes
\[
e^{\lambda p}e^{2\lambda}E_\lambda(p+2)E_\lambda(p-2).
\]

It remains to distinguish the physical height from the packet endpoint.
Lemma~\ref{lem:pre-HLZ-packet-separation} places the two packet legs in a common
output channel before scalar HLZ deformation.  Restricting that output to the
common packet prefix therefore replaces each full one-sided feature
$E_\lambda(q)$ by $E_\rho(q)$, and only those features.  The explicit powers
$X^a$ and $X^\delta$ are scalar archimedean multipliers and are untouched by
this restriction, exactly as in Proposition~\ref{prop:X-tagging}.  Applying
this replacement to the three full-fibre calculations above proves
\eqref{eq:independent-three-face-audit}.
\end{proof}

\begin{proposition}[Completed principal local profile]
\label{prop:completed-local-profile}
Let $\tau\asymp T$ be the physical height and let $\tau_b$ be the reference
height of its smooth block, of length $H=T/L^D$. Write
\[
\lambda(\tau):=\frac{\log(\tau/2\pi)}L,
\qquad
\mu_b:=\frac{2F_b}{L},
\]
where $F_b$ is the frozen value of the half-self factor $f$ on that block.
Stirling's formula and the mean-value theorem give, uniformly for
$|\tau-\tau_b|\ll H$,
\begin{equation}
\mu_b=\lambda(\tau)+O(L^{-D-1})+O((TL)^{-1}).
\label{eq:mu-lambda-close}
\end{equation}
Let $\rho$ be the pair-independent common $N$-prefix endpoint. Before taking
the normalized curvature coefficient, the completed model principal source has
local profile
\begin{equation}
\boxed{
\mathscr G_{\lambda(\tau),\mu_b,\rho}(p)
=e^{\lambda(\tau)p}\left[
\mu_b E_\rho(2p)-2E_\rho(p)^2
+e^{2\lambda(\tau)}E_\rho(p+2)E_\rho(p-2)
\right].
}
\label{eq:two-parameter-profile}
\end{equation}
Consequently the frozen profile $\lambda=\mu_b=1$ is
\begin{equation}
\boxed{
G(\rho)=[p^2]\,e^p\left[
E_\rho(2p)-2E_\rho(p)^2
+e^2E_\rho(p+2)E_\rho(p-2)
\right].
}
\label{eq:G-from-source}
\end{equation}
\end{proposition}

\begin{proof}
Use the completed per-orientation source
$H_{\rm mod}=F_b-2P+P(s-2/L)$ from \eqref{eq:Hmod}. By
Proposition~\ref{prop:trivial-outer-twist}, all auxiliary outer indices are
trivial, so the common-product coordinate reduces to $N=\ell$ and there is no
outward arithmetic annulus.
The common-prefix projection is taken before the local residue, so after
translating the normalized logarithmic carrier coordinate every principal fibre
integral is over the same interval $0\le t\le\rho$.

The self term is read after the two Hermitian orientations are reassembled.
Before freezing, the full $J$ source is $2f=-\chi'/\chi$; on the block under
consideration its principal frozen value is $2F_b$. The local pole extraction
has the same common positive $L^2$ Perron amplitude as the other principal
faces, so after this common amplitude is factored out the self coefficient is
exactly
\[
\frac{2F_b}{L}=\mu_b.
\]
Proposition~\ref{prop:J-full-fibre} gives the full glued fibre and
Proposition~\ref{prop:X-tagging} keeps the explicit Hardy shift factor at the
physical height. Hence the normalized full-$J$ contribution is
\[
\mu_b e^{\lambda(\tau)p}E_\rho(2p).
\]
The approximation \eqref{eq:mu-lambda-close} is used only in the later
comparison with the frozen model; no factor $2$ is absorbed into the definition
of \eqref{eq:two-parameter-profile}.
The two unshifted logarithmic-derivative faces lie on the common diagonal
\eqref{eq:hlz-diagonal}. Their two one-sided fibre integrals are both
$E_\rho(p)$; the coefficient $-2P$ therefore gives
\[
-2e^{\lambda(\tau)p}E_\rho(p)^2.
\]
Finally the translated source $P(s-2/L)$ contributes the internal shifts
$+2$ and $-2$ in the two diagonal fibre variables. The explicit
functional-equation factor is archimedean, hence
$X^{2/L}=e^{2\lambda(\tau)}$ by Proposition~\ref{prop:X-tagging}. Its
contribution is therefore
\[
e^{\lambda(\tau)p}e^{2\lambda(\tau)}
E_\rho(p+2)E_\rho(p-2).
\]
Adding the three contributions proves
\eqref{eq:two-parameter-profile}. Setting $\lambda=\mu_b=1$ and extracting the
coefficient of $p^2$ gives \eqref{eq:G-from-source}. By
Proposition~\ref{prop:trivial-outer-twist}, no auxiliary outer arithmetic label
is present in this invariant-source calculation.
\end{proof}

\begin{lemma}[The local projector contracts the internal Dirichlet fibre]
\label{lem:internal-principal-contraction}
Assume Hypothesis~\ref{hyp:packet-bilinear-HLZ}. Fix a packet pair. In the
invariant-source specialization the common shifted-zeta and packet analysis
maps are fixed before the internal HLP index is exposed; there is no separate
outer arithmetic map. After the common diagonal is imposed, every exact HLP
level has a radial Dirichlet series
\[
\mathcal D_k(z)=\sum_{M\ge1}a_{k,M}M^{-z}.
\]
For $k\le K_0$, Hypothesis~\ref{hyp:packet-bilinear-HLZ}(ii) places this entire
Dirichlet sum inside the common output operator before the local Perron
residue; in particular $M$ is not retained as an additional external Hilbert
coordinate of the low-level principal packet form. The actual packet Mellin
weight $W$ therefore enters the residue functional
\begin{equation}
\boxed{
\mathcal L_{\rm pr,\le K_0}^{\rm ex}[W]
=\sum_{0\le k\le K_0}\omega_{k,L}
\Res_{u=0}\bigl(\mathcal F_{k,L}(1+u;p)W(u)\bigr).
}
\label{eq:internal-principal-functional}
\end{equation}
Thus no vector indexed by $M$ remains in the exact low-level common principal
fibre after the curvature coefficients $q\le2$ are extracted at $p=0$.
Lemma~\ref{lem:cluster-levelwise} shows more: in the invariant-source
specialization those exact low-level curvature coefficients have no additional
translated Perron poles, \eqref{eq:no-exact-translated-poles}, and the difference between the original
absolute-convergence Perron representation and the common $u=0$ residue is an
arbitrary-log contour error on the same common output channel.  The levels
$k>K_0$ are not acted on by this projector and retain their direct $M$-carriers
as in Lemma~\ref{lem:HLP-high-edge}. Replacing the level factors by
their leading Laurent jets and extending the resulting \emph{model} sum over
all $k$ gives the model projector \eqref{eq:model-PP-projector} and the radial
profile of Proposition~\ref{prop:completed-local-profile}. The extension beyond
$K_0$ is $K_{\rm lead,\infty}^{\rm pkt}-K_{\rm lead,\le K_0}^{\rm pkt}$; the
signed correction required in the exact-minus-model decomposition is its
negative, namely \eqref{eq:model-truncation-tail}.
\end{lemma}

\begin{proof}
On the safe spectral right edge the internal Dirichlet sums and the geometric
HLP hierarchy are absolutely convergent.  The point requiring an operator
statement is the passage from this carrier-indexed representation to a common
packet output.  That passage is precisely Hypothesis~\ref{hyp:packet-bilinear-HLZ}(ii):
for $k\le K_0$ the complete $M$-sum is formed inside the common output operator,
with the packet Mellin weight represented on the same output channel.  Only
after that summation is the resulting meromorphic coefficient deformed.
Lemma~\ref{lem:cluster-levelwise} gives the scalar pole statement and
Hypothesis~\ref{hyp:packet-bilinear-HLZ}(iii) gives its carrier-preserving
operator lift.  The only crossed pole is $u=0$, so summing the common residues
over $k\le K_0$ gives exactly
\eqref{eq:internal-principal-functional}; the connector and left-contour terms
belong to the common contour-error class.  No Dirichlet series is invoked on
the left of $\Re z=1$.  The levels $k>K_0$ are excluded from the hypothesis and
remain on the direct right-edge representation, hence retain their physical
$M$-carriers.

For the model comparison, replace the level factors by their leading Laurent
jets. The resulting infinite model series is absolutely convergent and
Lemma~\ref{lem:leading-jet-resummation} gives the self, base, and translated
coefficients of Proposition~\ref{prop:completed-local-profile}, hence
$\mathscr G_{\lambda,\mu_b,\rho}$. Its portion beyond $K_0$ is the radial model
truncation error \eqref{eq:model-truncation-tail}; the exact high-level carriers
remain in the direct decomposition of Lemma~\ref{lem:HLP-high-edge}.
\end{proof}
